\batchmode %% Suppresses most terminal output.
\documentclass{article}
\usepackage{color}
\definecolor{boxshade}{gray}{0.85}
\setlength{\textwidth}{360pt}
\setlength{\textheight}{541pt}
\usepackage{latexsym}
\usepackage{ifthen}
% \usepackage{color}
%%%%%%%%%%%%%%%%%%%%%%%%%%%%%%%%%%%%%%%%%%%%%%%%%%%%%%%%%%%%%%%%%%%%%%%%%%%%%
% SWITCHES                                                                  %
%%%%%%%%%%%%%%%%%%%%%%%%%%%%%%%%%%%%%%%%%%%%%%%%%%%%%%%%%%%%%%%%%%%%%%%%%%%%%
\newboolean{shading} 
\setboolean{shading}{false}
\makeatletter
 %% this is needed only when inserted into the file, not when
 %% used as a package file.
%%%%%%%%%%%%%%%%%%%%%%%%%%%%%%%%%%%%%%%%%%%%%%%%%%%%%%%%%%%%%%%%%%%%%%%%%%%%%
%                                                                           %
% DEFINITIONS OF SYMBOL-PRODUCING COMMANDS                                  %
%                                                                           %
%    TLA+      LaTeX                                                        %
%    symbol    command                                                      %
%    ------    -------                                                      %
%    =>        \implies                                                     %
%    <:        \ltcolon                                                     %
%    :>        \colongt                                                     %
%    ==        \defeq                                                       %
%    ..        \dotdot                                                      %
%    ::        \coloncolon                                                  %
%    =|        \eqdash                                                      %
%    ++        \pp                                                          %
%    --        \mm                                                          %
%    **        \stst                                                        %
%    //        \slsl                                                        %
%    ^         \ct                                                          %
%    \A        \A                                                           %
%    \E        \E                                                           %
%    \AA       \AA                                                          %
%    \EE       \EE                                                          %
%%%%%%%%%%%%%%%%%%%%%%%%%%%%%%%%%%%%%%%%%%%%%%%%%%%%%%%%%%%%%%%%%%%%%%%%%%%%%
\newlength{\symlength}
\newcommand{\implies}{\Rightarrow}
\newcommand{\ltcolon}{\mathrel{<\!\!\mbox{:}}}
\newcommand{\colongt}{\mathrel{\!\mbox{:}\!\!>}}
\newcommand{\defeq}{\;\mathrel{\smash   %% keep this symbol from being too tall
    {{\stackrel{\scriptscriptstyle\Delta}{=}}}}\;}
\newcommand{\dotdot}{\mathrel{\ldotp\ldotp}}
\newcommand{\coloncolon}{\mathrel{::\;}}
\newcommand{\eqdash}{\mathrel = \joinrel \hspace{-.28em}|}
\newcommand{\pp}{\mathbin{++}}
\newcommand{\mm}{\mathbin{--}}
\newcommand{\stst}{*\!*}
\newcommand{\slsl}{/\!/}
\newcommand{\ct}{\hat{\hspace{.4em}}}
\newcommand{\A}{\forall}
\newcommand{\E}{\exists}
\renewcommand{\AA}{\makebox{$\raisebox{.05em}{\makebox[0pt][l]{%
   $\forall\hspace{-.517em}\forall\hspace{-.517em}\forall$}}%
   \forall\hspace{-.517em}\forall \hspace{-.517em}\forall\,$}}
\newcommand{\EE}{\makebox{$\raisebox{.05em}{\makebox[0pt][l]{%
   $\exists\hspace{-.517em}\exists\hspace{-.517em}\exists$}}%
   \exists\hspace{-.517em}\exists\hspace{-.517em}\exists\,$}}
\newcommand{\whileop}{\.{\stackrel
  {\mbox{\raisebox{-.3em}[0pt][0pt]{$\scriptscriptstyle+\;\,$}}}%
  {-\hspace{-.16em}\triangleright}}}

% Commands are defined to produce the upper-case keywords.
% Note that some have space after them.
\newcommand{\ASSUME}{\textsc{assume }}
\newcommand{\ASSUMPTION}{\textsc{assumption }}
\newcommand{\AXIOM}{\textsc{axiom }}
\newcommand{\BOOLEAN}{\textsc{boolean }}
\newcommand{\CASE}{\textsc{case }}
\newcommand{\CONSTANT}{\textsc{constant }}
\newcommand{\CONSTANTS}{\textsc{constants }}
\newcommand{\ELSE}{\settowidth{\symlength}{\THEN}%
   \makebox[\symlength][l]{\textsc{ else}}}
\newcommand{\EXCEPT}{\textsc{ except }}
\newcommand{\EXTENDS}{\textsc{extends }}
\newcommand{\FALSE}{\textsc{false}}
\newcommand{\IF}{\textsc{if }}
\newcommand{\IN}{\settowidth{\symlength}{\LET}%
   \makebox[\symlength][l]{\textsc{in}}}
\newcommand{\INSTANCE}{\textsc{instance }}
\newcommand{\LET}{\textsc{let }}
\newcommand{\LOCAL}{\textsc{local }}
\newcommand{\MODULE}{\textsc{module }}
\newcommand{\OTHER}{\textsc{other }}
\newcommand{\STRING}{\textsc{string}}
\newcommand{\THEN}{\textsc{ then }}
\newcommand{\THEOREM}{\textsc{theorem }}
\newcommand{\LEMMA}{\textsc{lemma }}
\newcommand{\PROPOSITION}{\textsc{proposition }}
\newcommand{\COROLLARY}{\textsc{corollary }}
\newcommand{\TRUE}{\textsc{true}}
\newcommand{\VARIABLE}{\textsc{variable }}
\newcommand{\VARIABLES}{\textsc{variables }}
\newcommand{\WITH}{\textsc{ with }}
\newcommand{\WF}{\textrm{WF}}
\newcommand{\SF}{\textrm{SF}}
\newcommand{\CHOOSE}{\textsc{choose }}
\newcommand{\ENABLED}{\textsc{enabled }}
\newcommand{\UNCHANGED}{\textsc{unchanged }}
\newcommand{\SUBSET}{\textsc{subset }}
\newcommand{\UNION}{\textsc{union }}
\newcommand{\DOMAIN}{\textsc{domain }}
% Added for tla2tex
\newcommand{\BY}{\textsc{by }}
\newcommand{\OBVIOUS}{\textsc{obvious }}
\newcommand{\HAVE}{\textsc{have }}
\newcommand{\QED}{\textsc{qed }}
\newcommand{\TAKE}{\textsc{take }}
\newcommand{\DEF}{\textsc{ def }}
\newcommand{\HIDE}{\textsc{hide }}
\newcommand{\RECURSIVE}{\textsc{recursive }}
\newcommand{\USE}{\textsc{use }}
\newcommand{\DEFINE}{\textsc{define }}
\newcommand{\PROOF}{\textsc{proof }}
\newcommand{\WITNESS}{\textsc{witness }}
\newcommand{\PICK}{\textsc{pick }}
\newcommand{\DEFS}{\textsc{defs }}
\newcommand{\PROVE}{\settowidth{\symlength}{\ASSUME}%
   \makebox[\symlength][l]{\textsc{prove}}\@s{-4.1}}%
  %% The \@s{-4.1) is a kludge added on 24 Oct 2009 [happy birthday, Ellen]
  %% so the correct alignment occurs if the user types
  %%   ASSUME X
  %%   PROVE  Y
  %% because it cancels the extra 4.1 pts added because of the 
  %% extra space after the PROVE.  This seems to works OK.
  %% However, the 4.1 equals Parameters.LaTeXLeftSpace(1) and
  %% should be changed if that method ever changes.
\newcommand{\SUFFICES}{\textsc{suffices }}
\newcommand{\NEW}{\textsc{new }}
\newcommand{\LAMBDA}{\textsc{lambda }}
\newcommand{\STATE}{\textsc{state }}
\newcommand{\ACTION}{\textsc{action }}
\newcommand{\TEMPORAL}{\textsc{temporal }}
\newcommand{\ONLY}{\textsc{only }}              %% added by LL on 2 Oct 2009
\newcommand{\OMITTED}{\textsc{omitted }}        %% added by LL on 31 Oct 2009
\newcommand{\@pfstepnum}[2]{\ensuremath{\langle#1\rangle}\textrm{#2}}
\newcommand{\bang}{\@s{1}\mbox{\small !}\@s{1}}
%% We should format || differently in PlusCal code than in TLA+ formulas.
\newcommand{\p@barbar}{\ifpcalsymbols
   \,\,\rule[-.25em]{.075em}{1em}\hspace*{.2em}\rule[-.25em]{.075em}{1em}\,\,%
   \else \,||\,\fi}
%% PlusCal keywords
\newcommand{\p@fair}{\textbf{fair }}
\newcommand{\p@semicolon}{\textbf{\,; }}
\newcommand{\p@algorithm}{\textbf{algorithm }}
\newcommand{\p@mmfair}{\textbf{-{}-fair }}
\newcommand{\p@mmalgorithm}{\textbf{-{}-algorithm }}
\newcommand{\p@assert}{\textbf{assert }}
\newcommand{\p@await}{\textbf{await }}
\newcommand{\p@begin}{\textbf{begin }}
\newcommand{\p@end}{\textbf{end }}
\newcommand{\p@call}{\textbf{call }}
\newcommand{\p@define}{\textbf{define }}
\newcommand{\p@do}{\textbf{ do }}
\newcommand{\p@either}{\textbf{either }}
\newcommand{\p@or}{\textbf{or }}
\newcommand{\p@goto}{\textbf{goto }}
\newcommand{\p@if}{\textbf{if }}
\newcommand{\p@then}{\,\,\textbf{then }}
\newcommand{\p@else}{\ifcsyntax \textbf{else } \else \,\,\textbf{else }\fi}
\newcommand{\p@elsif}{\,\,\textbf{elsif }}
\newcommand{\p@macro}{\textbf{macro }}
\newcommand{\p@print}{\textbf{print }}
\newcommand{\p@procedure}{\textbf{procedure }}
\newcommand{\p@process}{\textbf{process }}
\newcommand{\p@return}{\textbf{return}}
\newcommand{\p@skip}{\textbf{skip}}
\newcommand{\p@variable}{\textbf{variable }}
\newcommand{\p@variables}{\textbf{variables }}
\newcommand{\p@while}{\textbf{while }}
\newcommand{\p@when}{\textbf{when }}
\newcommand{\p@with}{\textbf{with }}
\newcommand{\p@lparen}{\textbf{(\,\,}}
\newcommand{\p@rparen}{\textbf{\,\,) }}   
\newcommand{\p@lbrace}{\textbf{\{\,\,}}   
\newcommand{\p@rbrace}{\textbf{\,\,\} }}

%%%%%%%%%%%%%%%%%%%%%%%%%%%%%%%%%%%%%%%%%%%%%%%%%%%%%%%%%
% REDEFINE STANDARD COMMANDS TO MAKE THEM FORMAT BETTER %
%                                                       %
% We redefine \in and \notin                            %
%%%%%%%%%%%%%%%%%%%%%%%%%%%%%%%%%%%%%%%%%%%%%%%%%%%%%%%%%
\renewcommand{\_}{\rule{.4em}{.06em}\hspace{.05em}}
\newlength{\equalswidth}
\let\oldin=\in
\let\oldnotin=\notin
\renewcommand{\in}{%
   {\settowidth{\equalswidth}{$\.{=}$}\makebox[\equalswidth][c]{$\oldin$}}}
\renewcommand{\notin}{%
   {\settowidth{\equalswidth}{$\.{=}$}\makebox[\equalswidth]{$\oldnotin$}}}


%%%%%%%%%%%%%%%%%%%%%%%%%%%%%%%%%%%%%%%%%%%%%%%%%%%%
%                                                  %
% HORIZONTAL BARS:                                 %
%                                                  %
%   \moduleLeftDash    |~~~~~~~~~~                 %
%   \moduleRightDash    ~~~~~~~~~~|                %
%   \midbar            |----------|                %
%   \bottombar         |__________|                %
%%%%%%%%%%%%%%%%%%%%%%%%%%%%%%%%%%%%%%%%%%%%%%%%%%%%
\newlength{\charwidth}\settowidth{\charwidth}{{\small\tt M}}
\newlength{\boxrulewd}\setlength{\boxrulewd}{.4pt}
\newlength{\boxlineht}\setlength{\boxlineht}{.5\baselineskip}
\newcommand{\boxsep}{\charwidth}
\newlength{\boxruleht}\setlength{\boxruleht}{.5ex}
\newlength{\boxruledp}\setlength{\boxruledp}{-\boxruleht}
\addtolength{\boxruledp}{\boxrulewd}
\newcommand{\boxrule}{\leaders\hrule height \boxruleht depth \boxruledp
                      \hfill\mbox{}}
\newcommand{\@computerule}{%
  \setlength{\boxruleht}{.5ex}%
  \setlength{\boxruledp}{-\boxruleht}%
  \addtolength{\boxruledp}{\boxrulewd}}

\newcommand{\bottombar}{\hspace{-\boxsep}%
  \raisebox{-\boxrulewd}[0pt][0pt]{\rule[.5ex]{\boxrulewd}{\boxlineht}}%
  \boxrule
  \raisebox{-\boxrulewd}[0pt][0pt]{%
      \rule[.5ex]{\boxrulewd}{\boxlineht}}\hspace{-\boxsep}\vspace{0pt}}

\newcommand{\moduleLeftDash}%
   {\hspace*{-\boxsep}%
     \raisebox{-\boxlineht}[0pt][0pt]{\rule[.5ex]{\boxrulewd
               }{\boxlineht}}%
    \boxrule\hspace*{.4em }}

\newcommand{\moduleRightDash}%
    {\hspace*{.4em}\boxrule
    \raisebox{-\boxlineht}[0pt][0pt]{\rule[.5ex]{\boxrulewd
               }{\boxlineht}}\hspace{-\boxsep}}%\vspace{.2em}

\newcommand{\midbar}{\hspace{-\boxsep}\raisebox{-.5\boxlineht}[0pt][0pt]{%
   \rule[.5ex]{\boxrulewd}{\boxlineht}}\boxrule\raisebox{-.5\boxlineht%
   }[0pt][0pt]{\rule[.5ex]{\boxrulewd}{\boxlineht}}\hspace{-\boxsep}}

%%%%%%%%%%%%%%%%%%%%%%%%%%%%%%%%%%%%%%%%%%%%%%%%%%%%%%%%%%%%%%%%%%%%%%%%%%%%%
% FORMATING COMMANDS                                                        %
%%%%%%%%%%%%%%%%%%%%%%%%%%%%%%%%%%%%%%%%%%%%%%%%%%%%%%%%%%%%%%%%%%%%%%%%%%%%%

%%%%%%%%%%%%%%%%%%%%%%%%%%%%%%%%%%%%%%%%%%%%%%%%%%%%%%%%%%%%%%%%%%%%%%%%%%%%%
% PLUSCAL SHADING                                                           %
%%%%%%%%%%%%%%%%%%%%%%%%%%%%%%%%%%%%%%%%%%%%%%%%%%%%%%%%%%%%%%%%%%%%%%%%%%%%%

% The TeX pcalshading switch is set on to cause PlusCal shading to be
% performed.  This changes the behavior of the following commands and
% environments to cause full-width shading to be performed on all lines.
% 
%   \tstrut \@x cpar mcom \@pvspace
% 
% The TeX pcalsymbols switch is turned on when typesetting a PlusCal algorithm,
% whether or not shading is being performed.  It causes symbols (other than
% parentheses and braces and PlusCal-only keywords) that should be typeset
% differently depending on whether they are in an algorithm to be typeset
% appropriately.  Currently, the only such symbol is "||".
%
% The TeX csyntax switch is turned on when typesetting a PlusCal algorithm in
% c-syntax.  This allows symbols to be format differently in the two syntaxes.
% The "else" keyword is the only one that is.

\newif\ifpcalshading \pcalshadingfalse
\newif\ifpcalsymbols \pcalsymbolsfalse
\newif\ifcsyntax     \csyntaxtrue

% The \@pvspace command makes a vertical space.  It uses \vspace
% except with \ifpcalshading, in which case it sets \pvcalvspace
% and the space is added by a following \@x command.
%
\newlength{\pcalvspace}\setlength{\pcalvspace}{0pt}%
\newcommand{\@pvspace}[1]{%
  \ifpcalshading
     \par\global\setlength{\pcalvspace}{#1}%
  \else
     \par\vspace{#1}%
  \fi
}

% The lcom environment was changed to set \lcomindent equal to
% the indentation it produces.  This length is used by the
% cpar environment to make shading extend for the full width
% of the line.  This assumes that lcom environments are not
% nested.  I hope TLATeX does not nest them.
%
\newlength{\lcomindent}%
\setlength{\lcomindent}{0pt}%

%\tstrut: A strut to produce inter-paragraph space in a comment.
%\rstrut: A strut to extend the bottom of a one-line comment so
%         there's no break in the shading between comments on 
%         successive lines.
\newcommand\tstrut%
  {\raisebox{\vshadelen}{\raisebox{-.25em}{\rule{0pt}{1.15em}}}%
   \global\setlength{\vshadelen}{0pt}}
\newcommand\rstrut{\raisebox{-.25em}{\rule{0pt}{1.15em}}%
 \global\setlength{\vshadelen}{0pt}}


% \.{op} formats operator op in math mode with empty boxes on either side.
% Used because TeX otherwise vary the amount of space it leaves around op.
\renewcommand{\.}[1]{\ensuremath{\mbox{}#1\mbox{}}}

% \@s{n} produces an n-point space
\newcommand{\@s}[1]{\hspace{#1pt}}           

% \@x{txt} starts a specification line in the beginning with txt
% in the final LaTeX source.
\newlength{\@xlen}
\newcommand\xtstrut%
  {\setlength{\@xlen}{1.05em}%
   \addtolength{\@xlen}{\pcalvspace}%
    \raisebox{\vshadelen}{\raisebox{-.25em}{\rule{0pt}{\@xlen}}}%
   \global\setlength{\vshadelen}{0pt}%
   \global\setlength{\pcalvspace}{0pt}}

\newcommand{\@x}[1]{\par
  \ifpcalshading
  \makebox[0pt][l]{\shadebox{\xtstrut\hspace*{\textwidth}}}%
  \fi
  \mbox{$\mbox{}#1\mbox{}$}}  

% \@xx{txt} continues a specification line with the text txt.
\newcommand{\@xx}[1]{\mbox{$\mbox{}#1\mbox{}$}}  

% \@y{cmt} produces a one-line comment.
\newcommand{\@y}[1]{\mbox{\footnotesize\hspace{.65em}%
  \ifthenelse{\boolean{shading}}{%
      \shadebox{#1\hspace{-\the\lastskip}\rstrut}}%
               {#1\hspace{-\the\lastskip}\rstrut}}}

% \@z{cmt} produces a zero-width one-line comment.
\newcommand{\@z}[1]{\makebox[0pt][l]{\footnotesize
  \ifthenelse{\boolean{shading}}{%
      \shadebox{#1\hspace{-\the\lastskip}\rstrut}}%
               {#1\hspace{-\the\lastskip}\rstrut}}}


% \@w{str} produces the TLA+ string "str".
\newcommand{\@w}[1]{\textsf{``{#1}''}}             


%%%%%%%%%%%%%%%%%%%%%%%%%%%%%%%%%%%%%%%%%%%%%%%%%%%%%%%%%%%%%%%%%%%%%%%%%%%%%
% SHADING                                                                   %
%%%%%%%%%%%%%%%%%%%%%%%%%%%%%%%%%%%%%%%%%%%%%%%%%%%%%%%%%%%%%%%%%%%%%%%%%%%%%
\def\graymargin{1}
  % The number of points of margin in the shaded box.

% \definecolor{boxshade}{gray}{.85}
% Defines the darkness of the shading: 1 = white, 0 = black
% Added by TLATeX only if needed.

% \shadebox{txt} puts txt in a shaded box.
\newlength{\templena}
\newlength{\templenb}
\newsavebox{\tempboxa}
\newcommand{\shadebox}[1]{{\setlength{\fboxsep}{\graymargin pt}%
     \savebox{\tempboxa}{#1}%
     \settoheight{\templena}{\usebox{\tempboxa}}%
     \settodepth{\templenb}{\usebox{\tempboxa}}%
     \hspace*{-\fboxsep}\raisebox{0pt}[\templena][\templenb]%
        {\colorbox{boxshade}{\usebox{\tempboxa}}}\hspace*{-\fboxsep}}}

% \vshade{n} makes an n-point inter-paragraph space, with
%  shading if the `shading' flag is true.
\newlength{\vshadelen}
\setlength{\vshadelen}{0pt}
\newcommand{\vshade}[1]{\ifthenelse{\boolean{shading}}%
   {\global\setlength{\vshadelen}{#1pt}}%
   {\vspace{#1pt}}}

\newlength{\boxwidth}
\newlength{\multicommentdepth}

%%%%%%%%%%%%%%%%%%%%%%%%%%%%%%%%%%%%%%%%%%%%%%%%%%%%%%%%%%%%%%%%%%%%%%%%%%%%%
% THE cpar ENVIRONMENT                                                      %
% ^^^^^^^^^^^^^^^^^^^^                                                      %
% The LaTeX input                                                           %
%                                                                           %
%   \begin{cpar}{pop}{nest}{isLabel}{d}{e}{arg6}                            %
%     XXXXXXXXXXXXXXX                                                       %
%     XXXXXXXXXXXXXXX                                                       %
%     XXXXXXXXXXXXXXX                                                       %
%   \end{cpar}                                                              %
%                                                                           %
% produces one of two possible results.  If isLabel is the letter "T",      %
% it produces the following, where [label] is the result of typesetting     %
% arg6 in an LR box, and d is is a number representing a distance in        %
% points.                                                                   %
%                                                                           %
%   prevailing |<-- d -->[label]<- e ->XXXXXXXXXXXXXXX                      %
%         left |                       XXXXXXXXXXXXXXX                      %
%       margin |                       XXXXXXXXXXXXXXX                      %
%                                                                           %
% If isLabel is the letter "F", then it produces                            %
%                                                                           %
%   prevailing |<-- d -->XXXXXXXXXXXXXXXXXXXXXXX                            %
%         left |         <- e ->XXXXXXXXXXXXXXXX                            %
%       margin |                XXXXXXXXXXXXXXXX                            %
%                                                                           %
% where d and e are numbers representing distances in points.               %
%                                                                           %
% The prevailing left margin is the one in effect before the most recent    %
% pop (argument 1) cpar environments with "T" as the nest argument, where   %
% pop is a number \geq 0.                                                   %
%                                                                           %
% If the nest argument is the letter "T", then the prevailing left          %
% margin is moved to the left of the second (and following) lines of        %
% X's.  Otherwise, the prevailing left margin is left unchanged.            %
%                                                                           %
% An \unnest{n} command moves the prevailing left margin to where it was    %
% before the most recent n cpar environments with "T" as the nesting        %
% argument.                                                                 %
%                                                                           %
% The environment leaves no vertical space above or below it, or between    %
% its paragraphs.  (TLATeX inserts the proper amount of vertical space.)    %
%%%%%%%%%%%%%%%%%%%%%%%%%%%%%%%%%%%%%%%%%%%%%%%%%%%%%%%%%%%%%%%%%%%%%%%%%%%%%

\newcounter{pardepth}
\setcounter{pardepth}{0}

% \setgmargin{txt} defines \gmarginN to be txt, where N is \roman{pardepth}.
% \thegmargin equals \gmarginN, where N is \roman{pardepth}.
\newcommand{\setgmargin}[1]{%
  \expandafter\xdef\csname gmargin\roman{pardepth}\endcsname{#1}}
\newcommand{\thegmargin}{\csname gmargin\roman{pardepth}\endcsname}
\newcommand{\gmargin}{0pt}

\newsavebox{\tempsbox}

\newlength{\@cparht}
\newlength{\@cpardp}
\newenvironment{cpar}[6]{%
  \addtocounter{pardepth}{-#1}%
  \ifthenelse{\boolean{shading}}{\par\begin{lrbox}{\tempsbox}%
                                 \begin{minipage}[t]{\linewidth}}{}%
  \begin{list}{}{%
     \edef\temp{\thegmargin}
     \ifthenelse{\equal{#3}{T}}%
       {\settowidth{\leftmargin}{\hspace{\temp}\footnotesize #6\hspace{#5pt}}%
        \addtolength{\leftmargin}{#4pt}}%
       {\setlength{\leftmargin}{#4pt}%
        \addtolength{\leftmargin}{#5pt}%
        \addtolength{\leftmargin}{\temp}%
        \setlength{\itemindent}{-#5pt}}%
      \ifthenelse{\equal{#2}{T}}{\addtocounter{pardepth}{1}%
                                 \setgmargin{\the\leftmargin}}{}%
      \setlength{\labelwidth}{0pt}%
      \setlength{\labelsep}{0pt}%
      \setlength{\itemindent}{-\leftmargin}%
      \setlength{\topsep}{0pt}%
      \setlength{\parsep}{0pt}%
      \setlength{\partopsep}{0pt}%
      \setlength{\parskip}{0pt}%
      \setlength{\itemsep}{0pt}
      \setlength{\itemindent}{#4pt}%
      \addtolength{\itemindent}{-\leftmargin}}%
   \ifthenelse{\equal{#3}{T}}%
      {\item[\tstrut\footnotesize \hspace{\temp}{#6}\hspace{#5pt}]
        }%
      {\item[\tstrut\hspace{\temp}]%
         }%
   \footnotesize}
 {\hspace{-\the\lastskip}\tstrut
 \end{list}%
  \ifthenelse{\boolean{shading}}%
          {\end{minipage}%
           \end{lrbox}%
           \ifpcalshading
             \setlength{\@cparht}{\ht\tempsbox}%
             \setlength{\@cpardp}{\dp\tempsbox}%
             \addtolength{\@cparht}{.15em}%
             \addtolength{\@cpardp}{.2em}%
             \addtolength{\@cparht}{\@cpardp}%
            % I don't know what's going on here.  I want to add a
            % \pcalvspace high shaded line, but I don't know how to
            % do it.  A little trial and error shows that the following
            % does a reasonable job approximating that, eliminating
            % the line if \pcalvspace is small.
            \addtolength{\@cparht}{\pcalvspace}%
             \ifdim \pcalvspace > .8em
               \addtolength{\pcalvspace}{-.2em}%
               \hspace*{-\lcomindent}%
               \shadebox{\rule{0pt}{\pcalvspace}\hspace*{\textwidth}}\par
               \global\setlength{\pcalvspace}{0pt}%
               \fi
             \hspace*{-\lcomindent}%
             \makebox[0pt][l]{\raisebox{-\@cpardp}[0pt][0pt]{%
                 \shadebox{\rule{0pt}{\@cparht}\hspace*{\textwidth}}}}%
             \hspace*{\lcomindent}\usebox{\tempsbox}%
             \par
           \else
             \shadebox{\usebox{\tempsbox}}\par
           \fi}%
           {}%
  }

%%%%%%%%%%%%%%%%%%%%%%%%%%%%%%%%%%%%%%%%%%%%%%%%%%%%%%%%%%%%%%%%%%%%%%%%%%%%%%
% THE ppar ENVIRONMENT                                                       %
% ^^^^^^^^^^^^^^^^^^^^                                                       %
% The environment                                                            %
%                                                                            %
%   \begin{ppar} ... \end{ppar}                                              %
%                                                                            %
% is equivalent to                                                           %
%                                                                            %
%   \begin{cpar}{0}{F}{F}{0}{0}{} ... \end{cpar}                             %
%                                                                            %
% The environment is put around each line of the output for a PlusCal        %
% algorithm.                                                                 %
%%%%%%%%%%%%%%%%%%%%%%%%%%%%%%%%%%%%%%%%%%%%%%%%%%%%%%%%%%%%%%%%%%%%%%%%%%%%%%
%\newenvironment{ppar}{%
%  \ifthenelse{\boolean{shading}}{\par\begin{lrbox}{\tempsbox}%
%                                 \begin{minipage}[t]{\linewidth}}{}%
%  \begin{list}{}{%
%     \edef\temp{\thegmargin}
%        \setlength{\leftmargin}{0pt}%
%        \addtolength{\leftmargin}{\temp}%
%        \setlength{\itemindent}{0pt}%
%      \setlength{\labelwidth}{0pt}%
%      \setlength{\labelsep}{0pt}%
%      \setlength{\itemindent}{-\leftmargin}%
%      \setlength{\topsep}{0pt}%
%      \setlength{\parsep}{0pt}%
%      \setlength{\partopsep}{0pt}%
%      \setlength{\parskip}{0pt}%
%      \setlength{\itemsep}{0pt}
%      \setlength{\itemindent}{0pt}%
%      \addtolength{\itemindent}{-\leftmargin}}%
%      \item[\tstrut\hspace{\temp}]}%
% {\hspace{-\the\lastskip}\tstrut
% \end{list}%
%  \ifthenelse{\boolean{shading}}{\end{minipage}  
%                                 \end{lrbox}%
%                                 \shadebox{\usebox{\tempsbox}}\par}{}%
%  }

 %%% TESTING
 \newcommand{\xtest}[1]{\par
 \makebox[0pt][l]{\shadebox{\xtstrut\hspace*{\textwidth}}}%
 \mbox{$\mbox{}#1\mbox{}$}} 

% \newcommand{\xxtest}[1]{\par
% \makebox[0pt][l]{\shadebox{\xtstrut{#1}\hspace*{\textwidth}}}%
% \mbox{$\mbox{}#1\mbox{}$}} 

%\newlength{\pcalvspace}
%\setlength{\pcalvspace}{0pt}
% \newlength{\xxtestlen}
% \setlength{\xxtestlen}{0pt}
% \newcommand\xtstrut%
%   {\setlength{\xxtestlen}{1.15em}%
%    \addtolength{\xxtestlen}{\pcalvspace}%
%     \raisebox{\vshadelen}{\raisebox{-.25em}{\rule{0pt}{\xxtestlen}}}%
%    \global\setlength{\vshadelen}{0pt}%
%    \global\setlength{\pcalvspace}{0pt}}
   
   %%%% TESTING
   
   %% The xcpar environment
   %%  Note: overloaded use of \pcalvspace for testing.
   %%
%   \newlength{\xcparht}%
%   \newlength{\xcpardp}%
   
%   \newenvironment{xcpar}[6]{%
%  \addtocounter{pardepth}{-#1}%
%  \ifthenelse{\boolean{shading}}{\par\begin{lrbox}{\tempsbox}%
%                                 \begin{minipage}[t]{\linewidth}}{}%
%  \begin{list}{}{%
%     \edef\temp{\thegmargin}%
%     \ifthenelse{\equal{#3}{T}}%
%       {\settowidth{\leftmargin}{\hspace{\temp}\footnotesize #6\hspace{#5pt}}%
%        \addtolength{\leftmargin}{#4pt}}%
%       {\setlength{\leftmargin}{#4pt}%
%        \addtolength{\leftmargin}{#5pt}%
%        \addtolength{\leftmargin}{\temp}%
%        \setlength{\itemindent}{-#5pt}}%
%      \ifthenelse{\equal{#2}{T}}{\addtocounter{pardepth}{1}%
%                                 \setgmargin{\the\leftmargin}}{}%
%      \setlength{\labelwidth}{0pt}%
%      \setlength{\labelsep}{0pt}%
%      \setlength{\itemindent}{-\leftmargin}%
%      \setlength{\topsep}{0pt}%
%      \setlength{\parsep}{0pt}%
%      \setlength{\partopsep}{0pt}%
%      \setlength{\parskip}{0pt}%
%      \setlength{\itemsep}{0pt}%
%      \setlength{\itemindent}{#4pt}%
%      \addtolength{\itemindent}{-\leftmargin}}%
%   \ifthenelse{\equal{#3}{T}}%
%      {\item[\xtstrut\footnotesize \hspace{\temp}{#6}\hspace{#5pt}]%
%        }%
%      {\item[\xtstrut\hspace{\temp}]%
%         }%
%   \footnotesize}
% {\hspace{-\the\lastskip}\tstrut
% \end{list}%
%  \ifthenelse{\boolean{shading}}{\end{minipage}  
%                                 \end{lrbox}%
%                                 \setlength{\xcparht}{\ht\tempsbox}%
%                                 \setlength{\xcpardp}{\dp\tempsbox}%
%                                 \addtolength{\xcparht}{.15em}%
%                                 \addtolength{\xcpardp}{.2em}%
%                                 \addtolength{\xcparht}{\xcpardp}%
%                                 \hspace*{-\lcomindent}%
%                                 \makebox[0pt][l]{\raisebox{-\xcpardp}[0pt][0pt]{%
%                                      \shadebox{\rule{0pt}{\xcparht}\hspace*{\textwidth}}}}%
%                                 \hspace*{\lcomindent}\usebox{\tempsbox}%
%                                 \par}{}%
%  }
%  
% \newlength{\xmcomlen}
%\newenvironment{xmcom}[1]{%
%  \setcounter{pardepth}{0}%
%  \hspace{.65em}%
%  \begin{lrbox}{\alignbox}\sloppypar%
%      \setboolean{shading}{false}%
%      \setlength{\boxwidth}{#1pt}%
%      \addtolength{\boxwidth}{-.65em}%
%      \begin{minipage}[t]{\boxwidth}\footnotesize
%      \parskip=0pt\relax}%
%       {\end{minipage}\end{lrbox}%
%       \setlength{\xmcomlen}{\textwidth}%
%       \addtolength{\xmcomlen}{-\wd\alignbox}%
%       \settodepth{\alignwidth}{\usebox{\alignbox}}%
%       \global\setlength{\multicommentdepth}{\alignwidth}%
%       \setlength{\boxwidth}{\alignwidth}%
%       \global\addtolength{\alignwidth}{-\maxdepth}%
%       \addtolength{\boxwidth}{.1em}%
%       \raisebox{0pt}[0pt][0pt]{%
%        \ifthenelse{\boolean{shading}}%
%          {\hspace*{-\xmcomlen}\shadebox{\rule[-\boxwidth]{0pt}{0pt}%
%                                 \hspace*{\xmcomlen}\usebox{\alignbox}}}%
%          {\usebox{\alignbox}}}%
%       \vspace*{\alignwidth}\pagebreak[0]\vspace{-\alignwidth}\par}
% % a multi-line comment, whose first argument is its width in points.
%  
   
%%%%%%%%%%%%%%%%%%%%%%%%%%%%%%%%%%%%%%%%%%%%%%%%%%%%%%%%%%%%%%%%%%%%%%%%%%%%%%
% THE lcom ENVIRONMENT                                                       %
% ^^^^^^^^^^^^^^^^^^^^                                                       %
% A multi-line comment with no text to its left is typeset in an lcom        % 
% environment, whose argument is a number representing the indentation       % 
% of the left margin, in points.  All the text of the comment should be      % 
% inside cpar environments.                                                  % 
%%%%%%%%%%%%%%%%%%%%%%%%%%%%%%%%%%%%%%%%%%%%%%%%%%%%%%%%%%%%%%%%%%%%%%%%%%%%%%
\newenvironment{lcom}[1]{%
  \setlength{\lcomindent}{#1pt} % Added for PlusCal handling.
  \par\vspace{.2em}%
  \sloppypar
  \setcounter{pardepth}{0}%
  \footnotesize
  \begin{list}{}{%
    \setlength{\leftmargin}{#1pt}
    \setlength{\labelwidth}{0pt}%
    \setlength{\labelsep}{0pt}%
    \setlength{\itemindent}{0pt}%
    \setlength{\topsep}{0pt}%
    \setlength{\parsep}{0pt}%
    \setlength{\partopsep}{0pt}%
    \setlength{\parskip}{0pt}}
    \item[]}%
  {\end{list}\vspace{.3em}\setlength{\lcomindent}{0pt}%
 }


%%%%%%%%%%%%%%%%%%%%%%%%%%%%%%%%%%%%%%%%%%%%%%%%%%%%%%%%%%%%%%%%%%%%%%%%%%%%%
% THE mcom ENVIRONMENT AND \mutivspace COMMAND                              %
% ^^^^^^^^^^^^^^^^^^^^^^^^^^^^^^^^^^^^^^^^^^^^                              %
%                                                                           %
% A part of the spec containing a right-comment of the form                 %
%                                                                           %
%      xxxx (*************)                                                 %
%      yyyy (* ccccccccc *)                                                 %
%      ...  (* ccccccccc *)                                                 %
%           (* ccccccccc *)                                                 %
%           (* ccccccccc *)                                                 %
%           (*************)                                                 %
%                                                                           %
% is typeset by                                                             %
%                                                                           %
%     XXXX \begin{mcom}{d}                                                  %
%            CCCC ... CCC                                                   %
%          \end{mcom}                                                       %
%     YYYY ...                                                              %
%     \multivspace{n}                                                       %
%                                                                           %
% where the number d is the width in points of the comment, n is the        %
% number of xxxx, yyyy, ...  lines to the left of the comment.              %
% All the text of the comment should be typeset in cpar environments.       %
%                                                                           %
% This puts the comment into a single box (so no page breaks can occur      %
% within it).  The entire box is shaded iff the shading flag is true.       %
%%%%%%%%%%%%%%%%%%%%%%%%%%%%%%%%%%%%%%%%%%%%%%%%%%%%%%%%%%%%%%%%%%%%%%%%%%%%%
\newlength{\xmcomlen}%
\newenvironment{mcom}[1]{%
  \setcounter{pardepth}{0}%
  \hspace{.65em}%
  \begin{lrbox}{\alignbox}\sloppypar%
      \setboolean{shading}{false}%
      \setlength{\boxwidth}{#1pt}%
      \addtolength{\boxwidth}{-.65em}%
      \begin{minipage}[t]{\boxwidth}\footnotesize
      \parskip=0pt\relax}%
       {\end{minipage}\end{lrbox}%
       \setlength{\xmcomlen}{\textwidth}%       % For PlusCal shading
       \addtolength{\xmcomlen}{-\wd\alignbox}%  % For PlusCal shading
       \settodepth{\alignwidth}{\usebox{\alignbox}}%
       \global\setlength{\multicommentdepth}{\alignwidth}%
       \setlength{\boxwidth}{\alignwidth}%      % For PlusCal shading
       \global\addtolength{\alignwidth}{-\maxdepth}%
       \addtolength{\boxwidth}{.1em}%           % For PlusCal shading
      \raisebox{0pt}[0pt][0pt]{%
        \ifthenelse{\boolean{shading}}%
          {\ifpcalshading
             \hspace*{-\xmcomlen}%
             \shadebox{\rule[-\boxwidth]{0pt}{0pt}\hspace*{\xmcomlen}%
                          \usebox{\alignbox}}%
           \else
             \shadebox{\usebox{\alignbox}}
           \fi
          }%
          {\usebox{\alignbox}}}%
       \vspace*{\alignwidth}\pagebreak[0]\vspace{-\alignwidth}\par}
 % a multi-line comment, whose first argument is its width in points.


% \multispace{n} produces the vertical space indicated by "|"s in 
% this situation
%   
%     xxxx (*************)
%     xxxx (* ccccccccc *)
%      |   (* ccccccccc *)
%      |   (* ccccccccc *)
%      |   (* ccccccccc *)
%      |   (*************)
%
% where n is the number of "xxxx" lines.
\newcommand{\multivspace}[1]{\addtolength{\multicommentdepth}{-#1\baselineskip}%
 \addtolength{\multicommentdepth}{1.2em}%
 \ifthenelse{\lengthtest{\multicommentdepth > 0pt}}%
    {\par\vspace{\multicommentdepth}\par}{}}

%\newenvironment{hpar}[2]{%
%  \begin{list}{}{\setlength{\leftmargin}{#1pt}%
%                 \addtolength{\leftmargin}{#2pt}%
%                 \setlength{\itemindent}{-#2pt}%
%                 \setlength{\topsep}{0pt}%
%                 \setlength{\parsep}{0pt}%
%                 \setlength{\partopsep}{0pt}%
%                 \setlength{\parskip}{0pt}%
%                 \addtolength{\labelsep}{0pt}}%
%  \item[]\footnotesize}{\end{list}}
%    %%%%%%%%%%%%%%%%%%%%%%%%%%%%%%%%%%%%%%%%%%%%%%%%%%%%%%%%%%%%%%%%%%%%%%%%
%    % Typesets a sequence of paragraphs like this:                         %
%    %                                                                      %
%    %      left |<-- d1 --> XXXXXXXXXXXXXXXXXXXXXXXX                       %
%    %    margin |           <- d2 -> XXXXXXXXXXXXXXX                       %
%    %           |                    XXXXXXXXXXXXXXX                       %
%    %           |                                                          %
%    %           |                    XXXXXXXXXXXXXXX                       %
%    %           |                    XXXXXXXXXXXXXXX                       %
%    %                                                                      %
%    % where d1 = #1pt and d2 = #2pt, but with no vspace between            %
%    % paragraphs.                                                          %
%    %%%%%%%%%%%%%%%%%%%%%%%%%%%%%%%%%%%%%%%%%%%%%%%%%%%%%%%%%%%%%%%%%%%%%%%%

%%%%%%%%%%%%%%%%%%%%%%%%%%%%%%%%%%%%%%%%%%%%%%%%%%%%%%%%%%%%%%%%%%%%%%
% Commands for repeated characters that produce dashes.              %
%%%%%%%%%%%%%%%%%%%%%%%%%%%%%%%%%%%%%%%%%%%%%%%%%%%%%%%%%%%%%%%%%%%%%%
% \raisedDash{wd}{ht}{thk} makes a horizontal line wd characters wide, 
% raised a distance ht ex's above the baseline, with a thickness of 
% thk em's.
\newcommand{\raisedDash}[3]{\raisebox{#2ex}{\setlength{\alignwidth}{.5em}%
  \rule{#1\alignwidth}{#3em}}}

% The following commands take a single argument n and produce the
% output for n repeated characters, as follows
%   \cdash:    -
%   \tdash:    ~
%   \ceqdash:  =
%   \usdash:   _
\newcommand{\cdash}[1]{\raisedDash{#1}{.5}{.04}}
\newcommand{\usdash}[1]{\raisedDash{#1}{0}{.04}}
\newcommand{\ceqdash}[1]{\raisedDash{#1}{.5}{.08}}
\newcommand{\tdash}[1]{\raisedDash{#1}{1}{.08}}

\newlength{\spacewidth}
\setlength{\spacewidth}{.2em}
\newcommand{\e}[1]{\hspace{#1\spacewidth}}
%% \e{i} produces space corresponding to i input spaces.


%% Alignment-file Commands

\newlength{\alignboxwidth}
\newlength{\alignwidth}
\newsavebox{\alignbox}

% \al{i}{j}{txt} is used in the alignment file to put "%{i}{j}{wd}"
% in the log file, where wd is the width of the line up to that point,
% and txt is the following text.
\newcommand{\al}[3]{%
  \typeout{\%{#1}{#2}{\the\alignwidth}}%
  \cl{#3}}

%% \cl{txt} continues a specification line in the alignment file
%% with text txt.
\newcommand{\cl}[1]{%
  \savebox{\alignbox}{\mbox{$\mbox{}#1\mbox{}$}}%
  \settowidth{\alignboxwidth}{\usebox{\alignbox}}%
  \addtolength{\alignwidth}{\alignboxwidth}%
  \usebox{\alignbox}}

% \fl{txt} in the alignment file begins a specification line that
% starts with the text txt.
\newcommand{\fl}[1]{%
  \par
  \savebox{\alignbox}{\mbox{$\mbox{}#1\mbox{}$}}%
  \settowidth{\alignwidth}{\usebox{\alignbox}}%
  \usebox{\alignbox}}



  
%%%%%%%%%%%%%%%%%%%%%%%%%%%%%%%%%%%%%%%%%%%%%%%%%%%%%%%%%%%%%%%%%%%%%%%%%%%%%
% Ordinarily, TeX typesets letters in math mode in a special math italic    %
% font.  This makes it typeset "it" to look like the product of the         %
% variables i and t, rather than like the word "it".  The following         %
% commands tell TeX to use an ordinary italic font instead.                 %
%%%%%%%%%%%%%%%%%%%%%%%%%%%%%%%%%%%%%%%%%%%%%%%%%%%%%%%%%%%%%%%%%%%%%%%%%%%%%
\ifx\documentclass\undefined
\else
  \DeclareSymbolFont{tlaitalics}{\encodingdefault}{cmr}{m}{it}
  \let\itfam\symtlaitalics
\fi

\makeatletter
\newcommand{\tlx@c}{\c@tlx@ctr\advance\c@tlx@ctr\@ne}
\newcounter{tlx@ctr}
\c@tlx@ctr=\itfam \multiply\c@tlx@ctr"100\relax \advance\c@tlx@ctr "7061\relax
\mathcode`a=\tlx@c \mathcode`b=\tlx@c \mathcode`c=\tlx@c \mathcode`d=\tlx@c
\mathcode`e=\tlx@c \mathcode`f=\tlx@c \mathcode`g=\tlx@c \mathcode`h=\tlx@c
\mathcode`i=\tlx@c \mathcode`j=\tlx@c \mathcode`k=\tlx@c \mathcode`l=\tlx@c
\mathcode`m=\tlx@c \mathcode`n=\tlx@c \mathcode`o=\tlx@c \mathcode`p=\tlx@c
\mathcode`q=\tlx@c \mathcode`r=\tlx@c \mathcode`s=\tlx@c \mathcode`t=\tlx@c
\mathcode`u=\tlx@c \mathcode`v=\tlx@c \mathcode`w=\tlx@c \mathcode`x=\tlx@c
\mathcode`y=\tlx@c \mathcode`z=\tlx@c
\c@tlx@ctr=\itfam \multiply\c@tlx@ctr"100\relax \advance\c@tlx@ctr "7041\relax
\mathcode`A=\tlx@c \mathcode`B=\tlx@c \mathcode`C=\tlx@c \mathcode`D=\tlx@c
\mathcode`E=\tlx@c \mathcode`F=\tlx@c \mathcode`G=\tlx@c \mathcode`H=\tlx@c
\mathcode`I=\tlx@c \mathcode`J=\tlx@c \mathcode`K=\tlx@c \mathcode`L=\tlx@c
\mathcode`M=\tlx@c \mathcode`N=\tlx@c \mathcode`O=\tlx@c \mathcode`P=\tlx@c
\mathcode`Q=\tlx@c \mathcode`R=\tlx@c \mathcode`S=\tlx@c \mathcode`T=\tlx@c
\mathcode`U=\tlx@c \mathcode`V=\tlx@c \mathcode`W=\tlx@c \mathcode`X=\tlx@c
\mathcode`Y=\tlx@c \mathcode`Z=\tlx@c
\makeatother

%%%%%%%%%%%%%%%%%%%%%%%%%%%%%%%%%%%%%%%%%%%%%%%%%%%%%%%%%%
%                THE describe ENVIRONMENT                %
%%%%%%%%%%%%%%%%%%%%%%%%%%%%%%%%%%%%%%%%%%%%%%%%%%%%%%%%%%
%
%
% It is like the description environment except it takes an argument
% ARG that should be the text of the widest label.  It adjusts the
% indentation so each item with label LABEL produces
%%      LABEL             blah blah blah
%%      <- width of ARG ->blah blah blah
%%                        blah blah blah
\newenvironment{describe}[1]%
   {\begin{list}{}{\settowidth{\labelwidth}{#1}%
            \setlength{\labelsep}{.5em}%
            \setlength{\leftmargin}{\labelwidth}% 
            \addtolength{\leftmargin}{\labelsep}%
            \addtolength{\leftmargin}{\parindent}%
            \def\makelabel##1{\rm ##1\hfill}}%
            \setlength{\topsep}{0pt}}%% 
                % Sets \topsep to 0 to reduce vertical space above
                % and below embedded displayed equations
   {\end{list}}

%   For tlatex.TeX
\usepackage{verbatim}
\makeatletter
\def\tla{\let\%\relax%
         \@bsphack
         \typeout{\%{\the\linewidth}}%
             \let\do\@makeother\dospecials\catcode`\^^M\active
             \let\verbatim@startline\relax
             \let\verbatim@addtoline\@gobble
             \let\verbatim@processline\relax
             \let\verbatim@finish\relax
             \verbatim@}
\let\endtla=\@esphack

\let\pcal=\tla
\let\endpcal=\endtla
\let\ppcal=\tla
\let\endppcal=\endtla

% The tlatex environment is used by TLATeX.TeX to typeset TLA+.
% TLATeX.TLA starts its files by writing a \tlatex command.  This
% command/environment sets \parindent to 0 and defines \% to its
% standard definition because the writing of the log files is messed up
% if \% is defined to be something else.  It also executes
% \@computerule to determine the dimensions for the TLA horizonatl
% bars.
\newenvironment{tlatex}{\@computerule%
                        \setlength{\parindent}{0pt}%
                       \makeatletter\chardef\%=`\%}{}


% The notla environment produces no output.  You can turn a 
% tla environment to a notla environment to prevent tlatex.TeX from
% re-formatting the environment.

\def\notla{\let\%\relax%
         \@bsphack
             \let\do\@makeother\dospecials\catcode`\^^M\active
             \let\verbatim@startline\relax
             \let\verbatim@addtoline\@gobble
             \let\verbatim@processline\relax
             \let\verbatim@finish\relax
             \verbatim@}
\let\endnotla=\@esphack

\let\nopcal=\notla
\let\endnopcal=\endnotla
\let\noppcal=\notla
\let\endnoppcal=\endnotla

%%%%%%%%%%%%%%%%%%%%%%%% end of tlatex.sty file %%%%%%%%%%%%%%%%%%%%%%% 
% last modified on Fri  3 August 2012 at 14:23:49 PST by lamport

\begin{document}
\tlatex
\setboolean{shading}{true}
 \@x{}\moduleLeftDash\@xx{ {\MODULE}
 FormalizedCorollary}\moduleRightDash\@xx{}%
\@x{ {\EXTENDS} Integers ,\, FiniteSets ,\, Sequences ,\, Utils}%
\@pvspace{8.0pt}%
\@x{}%
\@y{\@s{0}%
 {\par}\cdash{20}
}%
\@xx{}%
\@x{}%
\@y{\@s{0}%
 Segments
}%
\@xx{}%
\@x{}%
\@y{\@s{0}%
 {\par}\cdash{20}
}%
\@xx{}%
\@pvspace{8.0pt}%
\@x{ Segment ( a ,\, b ) \.{\defeq}}%
\@x{\@s{16.4} {\langle} a ,\, b {\rangle}}%
\@pvspace{8.0pt}%
\@x{ Seg \.{\defeq}}%
 \@x{\@s{15.19} \{ s \.{\in} Int \.{\times} Int \.{:} s [ 1 ] \.{\leq} s [ 2 ]
 \}}%
\@pvspace{8.0pt}%
\@x{ EmptySeg \.{\defeq} {\langle} {\rangle}}%
\@y{\@s{0}%
 the empty segment is not part of \ensuremath{Seg
}}%
\@xx{}%
\@pvspace{8.0pt}%
\@x{ b ( s ) \.{\defeq}}%
\@x{\@s{18.12} s [ 1 ]}%
\@pvspace{8.0pt}%
\@x{ e ( s ) \.{\defeq}}%
\@x{\@s{18.24} s [ 2 ]}%
\@pvspace{8.0pt}%
\@x{ SegShiftLeft ( s ) \.{\defeq}}%
\@x{\@s{16.4} Segment ( b ( s ) \.{-} 1 ,\, e ( s ) \.{-} 1 )}%
\@pvspace{8.0pt}%
\@x{ SegTruncateLeft ( s ) \.{\defeq}}%
\@x{\@s{16.4} {\IF} b ( s ) \.{=} e ( s )}%
\@x{\@s{16.4} \.{\THEN} EmptySeg}%
\@x{\@s{16.4} \.{\ELSE} Segment ( b ( s ) \.{+} 1 ,\, e ( s ) )}%
\@pvspace{8.0pt}%
\@x{ Precedes ( s1 ,\, s2 ) \.{\defeq}}%
\@x{\@s{16.4}}%
\@y{\@s{0}%
 returns \ensuremath{s1} ≺ \ensuremath{s2
}}%
\@xx{}%
\@x{\@s{16.4} \.{\land} b ( s1 )\@s{0.12} \.{<} b ( s2 )}%
\@x{\@s{16.4} \.{\land} e ( s1 ) \.{<} e ( s2 )}%
\@pvspace{8.0pt}%
\@x{ SegSubsetEq ( s1 ,\, s2 ) \.{\defeq}}%
\@x{\@s{16.4}}%
\@y{\@s{0}%
 checks if \ensuremath{s1} is a subset of \ensuremath{s2
}}%
\@xx{}%
\@x{\@s{16.4} \.{\land} b ( s2 )\@s{0.12} \.{\leq} b ( s1 )}%
\@x{\@s{16.4} \.{\land} e ( s1 ) \.{\leq} e ( s2 )}%
\@pvspace{8.0pt}%
\@x{}%
\@y{\@s{0}%
 {\par}\cdash{20}
}%
\@xx{}%
\@x{}%
\@y{\@s{0}%
 \ensuremath{Multisegments
}}%
\@xx{}%
\@x{}%
\@y{\@s{0}%
 {\par}\cdash{20}
}%
\@xx{}%
\@pvspace{8.0pt}%
\@x{ Multisegments \.{\defeq}}%
\@x{\@s{16.4} Seq ( Seg )}%
\@pvspace{16.0pt}%
\@x{ SubMultisegment ( N ,\, M ) \.{\defeq}}%
\@x{\@s{16.4} SubBag ( N ,\, M )}%
\@pvspace{16.0pt}%
\@x{ SameMultisegment ( M ,\, N ) \.{\defeq}}%
\@x{\@s{16.4} SameBag ( M ,\, N )}%
\@pvspace{16.0pt}%
\@x{ IndexSet ( M ) \.{\defeq}}%
\@x{\@s{16.4} 1 \.{\dotdot} Len ( M )}%
\@pvspace{8.0pt}%
\@x{ Delta ( M ) \.{\defeq}}%
\@x{\@s{16.4}}%
\@y{\@s{0}%
 \ensuremath{Delta(M)} is the function \ensuremath{i \.{\mapsto} Delta\_i
}}%
\@xx{}%
\@x{\@s{16.4} [ i \.{\in} IndexSet ( M ) \.{\mapsto} M [ i ] ]}%
\@pvspace{16.0pt}%
\@x{ IsLadder ( M ) \.{\defeq}}%
\@x{\@s{16.4} \.{\land} M \.{\neq} {\langle} {\rangle}\@s{192.95}}%
\@y{\@s{0}%
 a ladder is nonempty
}%
\@xx{}%
\@x{\@s{16.4} \.{\land} \E\, R \.{\in} Permutations ( IndexSet ( M ) ) \.{:}}%
\@x{\@s{31.61} \A\, i \.{\in} 1 \.{\dotdot} ( Len ( M ) \.{-} 1 ) \.{:}}%
 \@x{\@s{42.93} Precedes ( M [ R [ i \.{+} 1 ] ] ,\, M [ R [ i ] ]
 )\@s{73.79}}%
\@y{\@s{0}%
 \ensuremath{Delta\_\{R[i\.{+}1]\}} ≺ \ensuremath{Delta\_\{R[i]\}
}}%
\@xx{}%
\@pvspace{8.0pt}%
\@x{}%
\@y{\@s{0}%
 {\par}\cdash{20}
}%
\@xx{}%
\@x{}%
\@y{\@s{0}%
 Depth
}%
\@xx{}%
\@x{}%
\@y{\@s{0}%
 {\par}\cdash{20}
}%
\@xx{}%
\@pvspace{8.0pt}%
\@x{ Depth ( M ) \.{\defeq}}%
\@x{\@s{16.4} \.{\LET}}%
\@x{\@s{32.8} IsDepthChain ( i ,\, j ,\, Chain ) \.{\defeq}}%
\@x{\@s{49.19} \.{\land} Chain [ 1 ] \.{=} i\@s{137.06}}%
\@y{\@s{0}%
 the chain starts at index \ensuremath{i
}}%
\@xx{}%
\@x{\@s{49.19} \.{\land} \A\, r \.{\in} 1 \.{\dotdot} j \.{:}}%
\@x{\@s{64.41} Precedes (}%
\@x{\@s{80.81} Delta ( M ) [ Chain [ r ] ] ,\,}%
\@x{\@s{80.81} Delta ( M ) [ Chain [ r \.{+} 1 ] ]}%
\@x{\@s{64.41} )\@s{183.88}}%
\@y{\@s{0}%
 each segment precedes the next one
}%
\@xx{}%
\@x{\@s{16.4} \.{\IN}}%
\@x{\@s{32.8} [ i \.{\in} IndexSet ( M ) \.{\mapsto}}%
\@x{\@s{43.76} Max ( \{}%
\@x{\@s{67.57} j \.{\in} 0 \.{\dotdot} ( Len ( M ) \.{-} 1 ) \.{:}}%
 \@x{\@s{80.28} \E\, Chain \.{\in} [ 1 \.{\dotdot} ( j \.{+} 1 )
 \.{\rightarrow} IndexSet ( M ) ] \.{:}}%
\@x{\@s{91.60} IsDepthChain ( i ,\, j ,\, Chain )\@s{49.19}}%
\@y{\@s{0}%
 there exists a chain of length \ensuremath{j} from \ensuremath{i
}}%
\@xx{}%
\@x{\@s{43.76} \} )}%
\@x{\@s{32.8} ]}%
\@pvspace{8.0pt}%
\@x{ d ( M ) \.{\defeq}}%
\@x{\@s{24.73}}%
\@y{\@s{0}%
 returns the depth of the multisegment \ensuremath{M
}}%
\@xx{}%
\@x{\@s{24.73} {\IF} M \.{=} {\langle} {\rangle}}%
\@x{\@s{24.73} \.{\THEN} 0}%
\@x{\@s{24.73} \.{\ELSE} Max ( Range ( Depth ( M ) ) )}%
\@pvspace{8.0pt}%
\@x{}%
\@y{\@s{0}%
 {\par}\cdash{20}
}%
\@xx{}%
\@x{}%
\@y{\@s{0}%
 The Map \ensuremath{K
}}%
\@xx{}%
\@x{}%
\@y{\@s{0}%
 {\par}\cdash{20}
}%
\@xx{}%
\@pvspace{8.0pt}%
\@x{ DepthFiber ( M ,\, k ) \.{\defeq}}%
\@x{\@s{16.4} Fiber ( Depth ( M ) ,\, k )\@s{160.04}}%
\@y{\@s{0}%
 \ensuremath{d\_M\.{\ct}\{\.{-}1\}(k)
}}%
\@xx{}%
\@pvspace{16.0pt}%
\@x{ AdmissibleEnumeration ( m ,\, k ) \.{\defeq}}%
\@x{\@s{16.4} \.{\LET}}%
\@x{\@s{32.8} F \.{\defeq} DepthFiber ( m ,\, k )\@s{127.10}}%
\@y{\@s{0}%
 depth fiber at \ensuremath{k
}}%
\@xx{}%
\@pvspace{8.0pt}%
\@x{\@s{32.8} IsAdmissible ( P ) \.{\defeq}}%
 \@x{\@s{49.19} \A\, r \.{\in} 1 \.{\dotdot} ( Cardinality ( F ) \.{-} 1 )
 \.{:}}%
\@x{\@s{60.52} SegSubsetEq (}%
\@x{\@s{76.92} Delta ( m ) [ P [ r \.{+} 1 ] ] ,\,}%
\@x{\@s{76.92} Delta ( m ) [ P [ r ] ]}%
\@x{\@s{60.52} )\@s{200.01}}%
\@y{\@s{0}%
 \ensuremath{P} satisfies the admissibility condition
}%
\@xx{}%
\@x{\@s{16.4} \.{\IN}}%
\@x{\@s{32.8} {\CHOOSE} P \.{\in} Permutations ( F ) \.{:}}%
\@x{\@s{49.19} IsAdmissible ( P )\@s{144.66}}%
\@y{\@s{0}%
 choose any admissible enumeration
}%
\@xx{}%
\@pvspace{16.0pt}%
\@x{ jIndex ( M ,\, k ) \.{\defeq}}%
\@x{\@s{16.4} {\IF} k \.{\in} Range ( Depth ( M ) )}%
\@x{\@s{16.4} \.{\THEN} Last ( AdmissibleEnumeration ( M ,\, k ) )}%
\@x{\@s{16.4} \.{\ELSE} 0}%
\@pvspace{16.0pt}%
\@x{ IVee ( M ) \.{\defeq}}%
\@x{\@s{19.96} \.{\LET}}%
\@x{\@s{36.36} Cycles \.{\defeq}}%
\@x{\@s{52.76} \{ AdmissibleEnumeration ( M ,\, k ) \.{:}}%
\@x{\@s{70.06} k \.{\in} Range ( Depth ( M ) ) \}\@s{102.5}}%
\@y{\@s{0}%
 one cycle for each depth fiber
}%
\@xx{}%
\@x{\@s{19.96} \.{\IN}}%
\@x{\@s{36.36} PermutationFromCycles (}%
\@x{\@s{52.76} IndexSet ( M ) ,\,}%
\@x{\@s{52.76} Cycles}%
\@x{\@s{36.36} )\@s{222.89}}%
\@y{\@s{0}%
 permutation \ensuremath{i \.{\mapsto} i\.{\ct}vee} induced by these cycles
}%
\@xx{}%
\@pvspace{8.0pt}%
\@x{}%
\@y{\@s{0}%
 \ensuremath{2.1a
}}%
\@xx{}%
\@x{ SegmentTransform ( M ,\, i ) \.{\defeq}}%
\@x{\@s{16.4}}%
\@y{\@s{0}%
 this is ∆′\ensuremath{\_i
}}%
\@xx{}%
\@x{\@s{16.4} Segment (}%
\@x{\@s{32.8} b ( Delta ( M ) [ i ] ) ,\,}%
\@x{\@s{32.8} e ( Delta ( M ) [ IVee ( M ) [ i ] ] )}%
\@x{\@s{16.4} )}%
\@pvspace{8.0pt}%
\@x{}%
\@y{\@s{0}%
 \ensuremath{2.1b
}}%
\@xx{}%
\@x{ DistinguishedIndices ( M ) \.{\defeq}}%
\@x{\@s{16.4} \{ jIndex ( M ,\, k ) \.{:} k \.{\in} Range ( Depth ( M ) ) \}}%
\@pvspace{8.0pt}%
\@x{ RemainingIndices ( M ) \.{\defeq}}%
\@x{\@s{16.4}}%
\@y{\@s{0}%
 this is I''\ensuremath{\.{=}}I\ensuremath{\backslash}I\mbox{'}
}%
\@xx{}%
\@x{\@s{16.4} IndexSet ( M ) \.{\,\backslash\,} DistinguishedIndices ( M )}%
\@pvspace{8.0pt}%
\@x{}%
\@y{\@s{0}%
 \ensuremath{2.1c
}}%
\@xx{}%
\@x{ l ( M ) \.{\defeq}}%
\@x{\@s{22.18} [ r \.{\in} 1 \.{\dotdot} ( d ( M ) \.{+} 1 ) \.{\mapsto}}%
\@x{\@s{34.35} SegmentTransform ( M ,\, jIndex ( M ,\, r \.{-} 1 ) )}%
\@x{\@s{22.18} ]}%
\@pvspace{16.0pt}%
\@x{ DerivedMultisegment ( M ) \.{\defeq}}%
\@x{\@s{16.4} EnumerateValues (}%
\@x{\@s{32.8} [ i \.{\in} RemainingIndices ( M ) \.{\mapsto}}%
\@x{\@s{43.76} SegmentTransform ( M ,\, i ) ]}%
\@x{\@s{16.4} )\@s{242.85}}%
\@y{\@s{0}%
 \ensuremath{m}\mbox{'} \ensuremath{\.{=}} sum of
 \ensuremath{Delta}\mbox{'}\_i for \ensuremath{i} in I''
}%
\@xx{}%
\@pvspace{16.0pt}%
\@x{}%
\@y{\@s{0}%
 2.2
}%
\@xx{}%
\@x{ bTransform ( M ,\, i ) \.{\defeq}}%
\@x{\@s{16.4} b ( SegmentTransform ( M ,\, i ) )}%
\@pvspace{8.0pt}%
\@x{ eTransform ( M ,\, i ) \.{\defeq}}%
\@x{\@s{16.4} e ( SegmentTransform ( M ,\, i ) )}%
\@pvspace{8.0pt}%
\@x{ K ( M ) \.{\defeq}}%
\@x{\@s{27.73} {\langle} l ( M ) ,\, DerivedMultisegment ( M ) {\rangle}}%
\@pvspace{8.0pt}%
\@x{}%
\@y{\@s{0}%
 {\par}\cdash{20}
}%
\@xx{}%
\@x{}%
\@y{\@s{0}%
 The single \ensuremath{MW} step
}%
\@xx{}%
\@x{}%
\@y{\@s{0}%
 {\par}\cdash{20}
}%
\@xx{}%
\@pvspace{8.0pt}%
\@x{ min ( M ) \.{\defeq}}%
 \@x{\@s{21.52} Min ( \{ b ( Delta ( M ) [ i ] ) \.{:} i \.{\in} IndexSet ( M
 ) \} )}%
\@pvspace{16.0pt}%
\@x{}%
\@y{\@s{0}%
 \ensuremath{3.1a\.{-}3.1c
}}%
\@xx{}%
\@x{ LeadingSequence ( M ) \.{\defeq}}%
\@x{\@s{16.4} \.{\LET}}%
\@x{\@s{32.8} CandidateSequences \.{\defeq}}%
\@x{\@s{49.19} {\UNION} \{}%
\@x{\@s{65.6} [ 1 \.{\dotdot} k \.{\rightarrow} IndexSet ( M ) ] \.{:}}%
\@x{\@s{65.6} k \.{\in} 1 \.{\dotdot} Len ( M )}%
\@x{\@s{49.19} \}\@s{208.94}}%
\@y{\@s{0}%
 all possible nonempty index sequences of relevant length
}%
\@xx{}%
\@pvspace{8.0pt}%
\@x{\@s{32.8} FirstOK ( i ) \.{\defeq}}%
 \@x{\@s{49.19} \.{\land}\@s{10.51} b ( Delta ( M ) [ i ] ) \.{=} min ( M
 )\@s{78.02}}%
\@y{\@s{0}%
 \ensuremath{Delta\_i} begins at \ensuremath{min} \ensuremath{M
}}%
\@xx{}%
\@x{\@s{49.19} \.{\land}\@s{10.51} \A\, q \.{\in} IndexSet ( M ) \.{:}}%
\@x{\@s{74.93} b ( Delta ( M ) [ q ] ) \.{=} min ( M )}%
\@x{\@s{74.93} \.{\implies}}%
 \@x{\@s{74.93} e ( Delta ( M ) [ i ] ) \.{\leq} e ( Delta ( M ) [ q ]
 )\@s{43.36}}%
\@y{\@s{0}%
 \ensuremath{Delta\_i} has minimal ending among these segments
}%
\@xx{}%
\@pvspace{8.0pt}%
\@x{\@s{32.8} NextCandidate ( i ,\, q ) \.{\defeq}}%
 \@x{\@s{49.19} \.{\land} Precedes ( Delta ( M ) [ i ] ,\, Delta ( M ) [ q ]
 )\@s{46.30}}%
\@y{\@s{0}%
 \ensuremath{Delta\_i} ≺ \ensuremath{Delta\_q
}}%
\@xx{}%
 \@x{\@s{49.19} \.{\land} b ( Delta ( M ) [ q ] ) \.{=} b ( Delta ( M ) [ i ]
 ) \.{+} 1\@s{41.0}}%
\@y{\@s{0}%
 \ensuremath{Delta\_q} begins exactly one later
}%
\@xx{}%
\@pvspace{8.0pt}%
\@x{\@s{32.8} NextOK ( i ,\, next ) \.{\defeq}}%
\@x{\@s{49.19} \.{\land} NextCandidate ( i ,\, next )\@s{98.39}}%
\@y{\@s{0}%
 next is an eligible successor
}%
\@xx{}%
\@x{\@s{49.19} \.{\land} \A\, q \.{\in} IndexSet ( M ) \.{:}}%
\@x{\@s{64.41} NextCandidate ( i ,\, q )}%
\@x{\@s{64.41} \.{\implies}}%
 \@x{\@s{64.41} e ( Delta ( M ) [ next ] ) \.{\leq} e ( Delta ( M ) [ q ]
 )\@s{36.89}}%
\@y{\@s{0}%
 next has minimal ending among eligible successors
}%
\@xx{}%
\@pvspace{8.0pt}%
\@x{\@s{32.8} IsValid ( S ) \.{\defeq}}%
\@x{\@s{49.19} \.{\land}\@s{3.97} FirstOK ( S [ 1 ] )\@s{139.54}}%
\@y{\@s{0}%
 first index satisfies (\ensuremath{3.1a})
}%
\@xx{}%
 \@x{\@s{49.19} \.{\land}\@s{3.97} \A\, r \.{\in} 1 \.{\dotdot} ( Len ( S )
 \.{-} 1 ) \.{:}}%
\@x{\@s{68.38} NextOK ( S [ r ] ,\, S [ r \.{+} 1 ] )\@s{95.37}}%
\@y{\@s{0}%
 successive indices satisfy (\ensuremath{3.1b})
}%
\@xx{}%
 \@x{\@s{49.19} \.{\land}\@s{3.97} {\lnot} \E\, q \.{\in} IndexSet ( M )
 \.{:}}%
\@x{\@s{70.95} NextCandidate ( Last ( S ) ,\, q )\@s{72.09}}%
\@y{\@s{0}%
 no successor exists after the last index: (\ensuremath{3.1c})
}%
\@xx{}%
\@pvspace{8.0pt}%
\@x{\@s{16.4} \.{\IN}}%
\@x{\@s{32.8} {\CHOOSE} S \.{\in} CandidateSequences \.{:}}%
\@x{\@s{49.19} IsValid ( S )\@s{151.7}}%
\@y{\@s{0}%
 choose one complete leading sequence
}%
\@xx{}%
\@pvspace{16.0pt}%
\@x{ IStar ( M ) \.{\defeq}}%
\@x{\@s{16.4}}%
\@y{\@s{0}%
 I* \ensuremath{\.{=} \{i\_1,\, \.{\dots},\, i\_k\}
}}%
\@xx{}%
\@x{\@s{16.4} Range ( LeadingSequence ( M ) )\@s{150.21}}%
\@y{\@s{0}%
 image of the chosen leading sequence
}%
\@xx{}%
\@pvspace{16.0pt}%
\@x{ DeltaStar ( M ,\, i ) \.{\defeq}}%
\@x{\@s{16.4} {\IF} i \.{\in} IStar ( M )\@s{205.38}}%
\@y{\@s{0}%
 if \ensuremath{i} is one of the leading indices
}%
\@xx{}%
\@x{\@s{16.4} \.{\THEN} SegTruncateLeft ( Delta ( M ) [ i ] )\@s{114.80}}%
\@y{\@s{0}%
 \ensuremath{Delta\_i}* \ensuremath{\.{=} \.{-}Delta\_i
}}%
\@xx{}%
\@x{\@s{16.4} \.{\ELSE} Delta ( M ) [ i ]\@s{193.30}}%
\@y{\@s{0}%
 otherwise \ensuremath{Delta\_i}* \ensuremath{\.{=} Delta\_i
}}%
\@xx{}%
\@pvspace{16.0pt}%
\@x{}%
\@y{\@s{0}%
 3.2
}%
\@xx{}%
\@x{ MCross ( M ) \.{\defeq}}%
\@x{\@s{16.4} \.{\LET}}%
\@x{\@s{32.8} KeptIndices \.{\defeq}}%
\@x{\@s{49.19} \{ i \.{\in} IndexSet ( M ) \.{:}}%
\@x{\@s{62.38} DeltaStar ( M ,\, i ) \.{\neq} EmptySeg \}\@s{69.7}}%
\@y{\@s{0}%
 discard empty transformed segments
}%
\@xx{}%
\@x{\@s{16.4} \.{\IN}}%
\@x{\@s{32.8} EnumerateValues (}%
\@x{\@s{49.19} [ i \.{\in} KeptIndices \.{\mapsto}}%
\@x{\@s{60.16} DeltaStar ( M ,\, i ) ]}%
\@x{\@s{32.8} )\@s{228.24}}%
\@y{\@s{0}%
 \ensuremath{m\.{\ct}}cross \ensuremath{\.{=}} sum of the nonempty
 \ensuremath{Delta\_i}*
}%
\@xx{}%
\@pvspace{16.0pt}%
\@x{ DeltaCircle ( M ) \.{\defeq}}%
\@x{\@s{16.4} Segment (}%
\@x{\@s{32.8} min ( M ) ,\,\@s{212.74}}%
\@y{\@s{0}%
 left endpoint is \ensuremath{min} \ensuremath{m
}}%
\@xx{}%
 \@x{\@s{32.8} min ( M ) \.{+} Len ( LeadingSequence ( M ) ) \.{-}
 1\@s{73.80}}%
\@y{\@s{0}%
 right endpoint is \ensuremath{min} \ensuremath{m \.{+} k \.{-} 1
}}%
\@xx{}%
\@x{\@s{16.4} )}%
\@pvspace{16.0pt}%
\@x{ MW ( M ) \.{\defeq}}%
 \@x{\@s{35.28} {\langle} MCross ( M ) ,\, DeltaCircle ( M )
 {\rangle}\@s{127.1}}%
\@y{\@s{0}%
 \ensuremath{MW(m) \.{=}} (\ensuremath{m\.{\ct}}cross,
 \ensuremath{Delta\.{\ct}circ(m)})
}%
\@xx{}%
\@pvspace{16.0pt}%
\@x{}%
\@y{\@s{0}%
 {\par}\cdash{20}
}%
\@xx{}%
\@x{}%
\@y{\@s{0}%
 Lemma 1
}%
\@xx{}%
\@x{}%
\@y{\@s{0}%
 {\par}\cdash{20}
}%
\@xx{}%
\@pvspace{8.0pt}%
\@x{ {\THEOREM} Lemma1 \.{\defeq}}%
 \@x{\@s{16.4} \A\, m \.{\in} ( Multisegments \.{\,\backslash\,} \{ {\langle}
 {\rangle} \} ) \.{:}}%
\@x{\@s{27.72} min ( m ) \.{<} min ( l ( m ) )}%
\@x{\@s{27.72} \.{\implies}}%
 \@x{\@s{27.72} \.{\land} DerivedMultisegment ( m ) \.{\neq} {\langle}
 {\rangle}}%
\@x{\@s{27.72} \.{\land} MCross ( m ) \.{\neq} {\langle} {\rangle}}%
\@pvspace{8.0pt}%
\@x{ {\PROOF} {\OMITTED}}%
\@pvspace{16.0pt}%
\@x{}%
\@y{\@s{0}%
 {\par}\cdash{20}
}%
\@xx{}%
\@x{}%
\@y{\@s{0}%
 The Statement - Corollary 3.4
}%
\@xx{}%
\@x{}%
\@y{\@s{0}%
 First alternative formulation
}%
\@xx{}%
\@x{}%
\@y{\@s{0}%
 {\par}\cdash{20}
}%
\@xx{}%
\@pvspace{8.0pt}%
\@x{ {\THEOREM} Corollary1 \.{\defeq}}%
 \@x{\@s{16.4} \A\, M \.{\in} ( Multisegments \.{\,\backslash\,} \{ {\langle}
 {\rangle} \} ) \.{:}\@s{112.48}}%
\@y{\@s{0}%
 for every nonempty multisegment \ensuremath{M
}}%
\@xx{}%
\@x{\@s{27.72} min ( M ) \.{<} min ( l ( M ) )\@s{140.50}}%
\@y{\@s{0}%
 assuming \ensuremath{min} \ensuremath{M \.{<} min} \ensuremath{l(M)
}}%
\@xx{}%
\@x{\@s{27.72} \.{\implies}}%
\@x{\@s{27.72} \.{\land} l ( M ) \.{=} l ( MCross ( M ) )\@s{127.13}}%
\@y{\@s{0}%
 \ensuremath{l(m) \.{=} l(m\.{\ct}cross)
}}%
\@xx{}%
\@x{\@s{27.72} \.{\land} DeltaCircle ( M )}%
\@x{\@s{42.93} \.{=} DeltaCircle ( DerivedMultisegment ( M ) )\@s{41.0}}%
\@y{\@s{0}%
 \ensuremath{Delta\.{\ct}circle(m) \.{=} Delta\.{\ct}circle(m\.{'})
}}%
\@xx{}%
\@x{\@s{27.72} \.{\land} SameMultisegment (}%
\@x{\@s{55.23} DerivedMultisegment ( MCross ( M ) ) ,\,}%
\@x{\@s{55.23} MCross ( DerivedMultisegment ( M ) )}%
\@x{\@s{38.83} )\@s{230.26}}%
\@y{\@s{0}%
 (\ensuremath{m\.{\ct}}cross)\mbox{'} \ensuremath{\.{=}}
 (\ensuremath{m}\mbox{'})\ensuremath{\.{\ct}}cross as multisegments
}%
\@xx{}%
\@pvspace{8.0pt}%
\@x{ {\PROOF} {\OMITTED}}%
\@pvspace{16.0pt}%
\@x{}%
\@y{\@s{0}%
 {\par}\cdash{20}
}%
\@xx{}%
\@x{}%
\@y{\@s{0}%
 The Statement - Corollary 3.4
}%
\@xx{}%
\@x{}%
\@y{\@s{0}%
 Second alternative formulation
}%
\@xx{}%
\@x{}%
\@y{\@s{0}%
 {\par}\cdash{20}
}%
\@xx{}%
\@pvspace{8.0pt}%
\@x{ {\THEOREM} Corollary2 \.{\defeq}}%
\@x{\@s{16.4} \.{\LET}}%
\@x{\@s{32.8} KxId ( P ) \.{\defeq}}%
\@x{\@s{55.12} \.{\LET}}%
\@x{\@s{71.52} KP \.{\defeq} K ( P [ 1 ] )\@s{131.78}}%
\@y{\@s{0}%
 apply \ensuremath{K} to the first component
}%
\@xx{}%
\@x{\@s{55.12} \.{\IN}}%
 \@x{\@s{71.52} {\langle} KP [ 1 ] ,\, KP [ 2 ] ,\, P [ 2 ]
 {\rangle}\@s{82.0}}%
\@y{\@s{0}%
 (\ensuremath{K} x \ensuremath{Id)(P)
}}%
\@xx{}%
\@pvspace{8.0pt}%
\@x{\@s{32.8} IdxMW ( P ) \.{\defeq}}%
\@x{\@s{49.19} \.{\LET}}%
\@x{\@s{65.6} MWP \.{\defeq} MW ( P [ 2 ] )\@s{114.8}}%
\@y{\@s{0}%
 apply \ensuremath{MW} to the second component
}%
\@xx{}%
\@x{\@s{49.19} \.{\IN}}%
 \@x{\@s{65.6} {\langle} P [ 1 ] ,\, MWP [ 1 ] ,\, MWP [ 2 ]
 {\rangle}\@s{76.85}}%
\@y{\@s{0}%
 (Id x \ensuremath{MW)(P)
}}%
\@xx{}%
\@pvspace{8.0pt}%
\@x{\@s{32.8} SameResult ( P ,\, Q ) \.{\defeq}}%
\@x{\@s{49.19} \.{\land} P [ 1 ] \.{=} Q [ 1 ]\@s{144.29}}%
\@y{\@s{0}%
 ladder components are equal
}%
\@xx{}%
\@x{\@s{49.19} \.{\land} SameMultisegment ( P [ 2 ] ,\, Q [ 2 ] )\@s{61.50}}%
\@y{\@s{0}%
 multisegment components are equal
}%
\@xx{}%
\@x{\@s{49.19} \.{\land} P [ 3 ] \.{=} Q [ 3 ]\@s{144.29}}%
\@y{\@s{0}%
 distinguished segments are equal
}%
\@xx{}%
\@pvspace{8.0pt}%
\@x{\@s{16.4} \.{\IN}}%
 \@x{\@s{32.8} \A\, M \.{\in} ( Multisegments \.{\,\backslash\,} \{ {\langle}
 {\rangle} \} ) \.{:}}%
\@x{\@s{44.12} min ( M ) \.{<} min ( l ( M ) )}%
\@x{\@s{44.12} \.{\implies}}%
\@x{\@s{44.12} SameResult (}%
\@x{\@s{60.52} KxId ( MW ( M ) ) ,\,}%
\@x{\@s{60.52} IdxMW ( K ( M ) )}%
\@x{\@s{44.12} )\@s{216.94}}%
\@y{\@s{0}%
 (\ensuremath{K} x \ensuremath{Id)(MW(m)) \.{=}} (\ensuremath{Id} x
 \ensuremath{MW)(K(m))
}}%
\@xx{}%
\@pvspace{8.0pt}%
\@x{ {\PROOF} {\OMITTED}}%
\@pvspace{16.0pt}%
\@x{}%
\@y{\@s{0}%
 {\par}\cdash{20}
}%
\@xx{}%
\@x{}%
\@y{\@s{0}%
 The Statement - Corollary 3.4
}%
\@xx{}%
\@x{}%
\@y{\@s{0}%
 Third alternative formulation
}%
\@xx{}%
\@x{}%
\@y{\@s{0}%
 {\par}\cdash{20}
}%
\@xx{}%
\@pvspace{8.0pt}%
\@x{ {\THEOREM} Corollary3 \.{\defeq}}%
 \@x{\@s{16.4} \A\, M \.{\in} ( Multisegments \.{\,\backslash\,} \{ {\langle}
 {\rangle} \} ) \.{:}}%
\@x{\@s{27.72} min ( M ) \.{<} min ( l ( M ) )}%
\@x{\@s{27.72} \.{\implies}}%
 \@x{\@s{27.72} \.{\land} K ( M ) [ 1 ] \.{=} K ( MW ( M ) [ 1 ] ) [ 1
 ]\@s{101.84}}%
\@y{\@s{0}%
 \ensuremath{l(m) \.{=} l(m\.{\ct}cross)
}}%
\@xx{}%
 \@x{\@s{27.72} \.{\land} MW ( M ) [ 2 ] \.{=} MW ( K ( M ) [ 2 ] ) [ 2
 ]\@s{90.2}}%
\@y{\@s{0}%
 \ensuremath{Delta\.{\ct}circ(m) \.{=} Delta\.{\ct}circ(m\.{'})
}}%
\@xx{}%
\@x{\@s{27.72} \.{\land} SameMultisegment (}%
\@x{\@s{55.23} K ( MW ( M ) [ 1 ] ) [ 2 ] ,\,}%
\@x{\@s{55.23} MW ( K ( M ) [ 2 ] ) [ 1 ]}%
\@x{\@s{38.83} )\@s{222.23}}%
\@y{\@s{0}%
 (\ensuremath{m\.{\ct}}cross)\mbox{'} \ensuremath{\.{=}}
 (\ensuremath{m}\mbox{'})\ensuremath{\.{\ct}}cross as multisegments
}%
\@xx{}%
\@pvspace{8.0pt}%
\@x{ {\PROOF} {\OMITTED}}%
\@pvspace{24.0pt}%
\@x{}\bottombar\@xx{}%
\setboolean{shading}{false}
\begin{lcom}{0}%
\begin{cpar}{0}{F}{F}{0}{0}{}%
Glossary
\end{cpar}%
\begin{cpar}{0}{F}{F}{0}{0}{}%
\ceqdash{76}
\end{cpar}%
\vshade{5.0}%
\begin{cpar}{0}{F}{F}{0}{0}{}%
Segment
 constructs the sequence \ensuremath{{\langle}a,\, b{\rangle}} (keeping in
 mind that a sequence is a
 function from indexes to values of the sequence).
 why choose a sequence and not a set\.{?}
 because what we mainly use and care about in the proof are the endpoints.
\end{cpar}%
\vshade{5.0}%
\begin{cpar}{0}{F}{F}{0}{0}{}%
\ensuremath{Seg
}%
\end{cpar}%
\begin{cpar}{0}{F}{F}{0}{0}{}%
 the set of all valid segments [a,\ensuremath{b}] with \ensuremath{a \.{\leq}
 b}.
\end{cpar}%
\vshade{5.0}%
\begin{cpar}{0}{F}{F}{0}{0}{}%
\ensuremath{EmptySeg
}%
\end{cpar}%
\begin{cpar}{0}{F}{F}{0}{0}{}%
 the empty sequence \ensuremath{{\langle}{\rangle}}, which we use to represent
 the empty segment.
 \ensuremath{EmptySeg} is not part of \ensuremath{Seg}.
\end{cpar}%
\vshade{5.0}%
\begin{cpar}{0}{F}{F}{0}{0}{}%
\ensuremath{b(s)
}%
\end{cpar}%
\begin{cpar}{0}{F}{F}{0}{0}{}%
returns the beginning/left endpoint of a segment \ensuremath{s}.
 if \ensuremath{s \.{=} {\langle}a,\,b{\rangle}}, then
 \ensuremath{b(s)\.{=}a}.
\end{cpar}%
\vshade{5.0}%
\begin{cpar}{0}{F}{F}{0}{0}{}%
\ensuremath{e(s)
}%
\end{cpar}%
\begin{cpar}{0}{F}{F}{0}{0}{}%
returns the ending/right endpoint of a segment \ensuremath{s}.
 if \ensuremath{s \.{=} {\langle}a,\,b{\rangle}}, then
 \ensuremath{e(s)\.{=}b}.
\end{cpar}%
\vshade{5.0}%
\begin{cpar}{0}{F}{F}{0}{0}{}%
\ensuremath{SegShiftLeft(s)
}%
\end{cpar}%
\begin{cpar}{0}{F}{F}{0}{0}{}%
shifts both endpoints of a segment one place to the left.
 [a,\ensuremath{b}] becomes [\ensuremath{a\.{-}1},\ensuremath{b\.{-}1}].
\end{cpar}%
\vshade{5.0}%
\begin{cpar}{0}{F}{F}{0}{0}{}%
\ensuremath{SegTruncateLeft(s)
}%
\end{cpar}%
\begin{cpar}{0}{F}{F}{0}{0}{}%
removes the left endpoint of a segment.
 if the segment is [a,\ensuremath{b}] with \ensuremath{a\.{<}b}, it returns
 [\ensuremath{a\.{+}1},\ensuremath{b}].
 if the segment is a singleton [a,a], it returns \ensuremath{EmptySeg}.
\end{cpar}%
\vshade{5.0}%
\begin{cpar}{0}{F}{F}{0}{0}{}%
\ensuremath{Precedes(s1,\,s2)
}%
\end{cpar}%
\begin{cpar}{0}{F}{F}{0}{0}{}%
checks whether \ensuremath{s1} ≺ \ensuremath{s2}.
 this means that both the beginning and the ending of \ensuremath{s1} are
 strictly
 smaller than those of \ensuremath{s2}.
\end{cpar}%
\vshade{5.0}%
\begin{cpar}{0}{F}{F}{0}{0}{}%
\ensuremath{SegSubsetEq(s1,\,s2)
}%
\end{cpar}%
\begin{cpar}{0}{F}{F}{0}{0}{}%
 checks whether \ensuremath{s1} is contained in \ensuremath{s2}, including
 equality.
 for \ensuremath{s1\.{=}[a1,\,b1]} and \ensuremath{s2\.{=}[a2,\,b2]}, this
 means
 \ensuremath{a2 \.{\leq} a1} and \ensuremath{b1 \.{\leq} b2}.
\end{cpar}%
\vshade{10.0}%
\begin{cpar}{0}{F}{F}{0}{0}{}%
\ensuremath{Multisegments
}%
\end{cpar}%
\begin{cpar}{0}{F}{F}{0}{0}{}%
the set of all multisegments, which are sequences of segments.
 why choose sequences of segments and not multisets/bags of segments\.{?}
 because we need to easily be able to use the indexes of the segments,
 even with repetition.
\end{cpar}%
\vshade{5.0}%
\begin{cpar}{0}{F}{F}{0}{0}{}%
this means that two different sequences may represent the same
 mathematical multisegment. we deal with this using
 \ensuremath{SameMultisegment
} rather than ordinary sequence equality.
\end{cpar}%
\vshade{5.0}%
\begin{cpar}{0}{F}{F}{0}{0}{}%
\ensuremath{SubMultisegment(N,\,M)
}%
\end{cpar}%
\begin{cpar}{0}{F}{F}{0}{0}{}%
checks whether \ensuremath{N} is a submultisegment of \ensuremath{M}.
 internally, this is \ensuremath{SubBag(N,\,M)}, meaning that every segment
 occurring in \ensuremath{N
} occurs in \ensuremath{M} at least as many times.
\end{cpar}%
\vshade{5.0}%
\begin{cpar}{0}{F}{F}{0}{0}{}%
\ensuremath{SameMultisegment(M,\,N)
}%
\end{cpar}%
\begin{cpar}{0}{F}{F}{0}{0}{}%
 checks whether \ensuremath{M} and \ensuremath{N} represent the same
 mathematical multisegment,
 even if they are different sequences.
 internally, this is \ensuremath{SameBag(M,\,N)}, so equality is determined
 by the
 segments and their multiplicities and not by their order.
\end{cpar}%
\vshade{5.0}%
\begin{cpar}{0}{F}{F}{0}{0}{}%
\ensuremath{IndexSet(M)
}%
\end{cpar}%
\begin{cpar}{0}{F}{F}{0}{0}{}%
 returns the set \ensuremath{1\.{\dotdot}Len(M)}, \ensuremath{i.e}. the set of
 global indexes of the segments
 in \ensuremath{M}.
\end{cpar}%
\vshade{5.0}%
\begin{cpar}{0}{F}{F}{0}{0}{}%
\ensuremath{Delta(M)
}%
\end{cpar}%
\begin{cpar}{0}{F}{F}{0}{0}{}%
 returns the function \ensuremath{i \.{\mapsto} Delta\_i}, where
 \ensuremath{Delta\_i} is the segment at
 global index \ensuremath{i}.
 since \ensuremath{M} itself is a sequence, this is extensionally the same
 function
 as \ensuremath{M}, but we keep \ensuremath{Delta(M)} in order to stay close
 to the notation
 of the paper.
\end{cpar}%
\vshade{5.0}%
\begin{cpar}{0}{F}{F}{0}{0}{}%
\ensuremath{IsLadder(M)
}%
\end{cpar}%
\begin{cpar}{0}{F}{F}{0}{0}{}%
checks whether \ensuremath{M} is a ladder.
 first, \ensuremath{M} must be nonempty.
 then there must exist a permutation \ensuremath{R} of
 \ensuremath{IndexSet(M)} such that the
 segments are in standard form, \ensuremath{i.e}.
 \ensuremath{Delta\_\{R[i\.{+}1]\}} ≺ \ensuremath{Delta\_\{R[i]\}} for every
 successive pair.
 using \ensuremath{Permutations} here already makes sure that \ensuremath{R}
 contains every
 global index exactly once.
\end{cpar}%
\vshade{10.0}%
\begin{cpar}{0}{F}{F}{0}{0}{}%
\ensuremath{Depth(M)
}%
\end{cpar}%
\begin{cpar}{0}{F}{F}{0}{0}{}%
 returns the function \ensuremath{i \.{\mapsto} d\_M(i)}, where
 \ensuremath{d\_M(i)} is the depth of the
 segment at index \ensuremath{i}.
 it does this by taking the maximum \ensuremath{j} such that there exists a
 depth chain
 of length \ensuremath{j} starting at index \ensuremath{i}.
\end{cpar}%
\vshade{5.0}%
\begin{cpar}{0}{F}{F}{0}{0}{}%
 internally, \ensuremath{IsDepthChain(i,\,j,\,Chain)} checks that
 \ensuremath{Chain[1]\.{=}i} and that
 for every \ensuremath{r} in [\ensuremath{j}],
 \ensuremath{Delta\_\{Chain[r]\}} ≺ \ensuremath{Delta\_\{Chain[r\.{+}1]\}}.
 we do not define \ensuremath{IsDepthChain} as an outer operator because it
 is only
 auxiliary machinery needed for defining \ensuremath{Depth}.
\end{cpar}%
\vshade{5.0}%
\begin{cpar}{0}{F}{F}{0}{0}{}%
\ensuremath{d(M)
}%
\end{cpar}%
\begin{cpar}{0}{F}{F}{0}{0}{}%
returns the depth of the multisegment \ensuremath{M}.
 if \ensuremath{M} is empty it returns 0.
 otherwise, it takes the maximum value in \ensuremath{Range(Depth(M))}.
\end{cpar}%
\vshade{10.0}%
\begin{cpar}{0}{F}{F}{0}{0}{}%
\ensuremath{DepthFiber(M,\,k)
}%
\end{cpar}%
\begin{cpar}{0}{F}{F}{0}{0}{}%
 returns the depth fiber \ensuremath{d\_M\.{\ct}(\.{-}1)(k)},
 \ensuremath{i.e}. the set of indexes \ensuremath{i} such that
 \ensuremath{Depth(M)[i]\.{=}k}.
 this is implemented using the general \ensuremath{Fiber} operator.
\end{cpar}%
\vshade{5.0}%
\begin{cpar}{0}{F}{F}{0}{0}{}%
\ensuremath{AdmissibleEnumeration(m,\,k)
}%
\end{cpar}%
\begin{cpar}{0}{F}{F}{0}{0}{}%
chooses an admissible enumeration of the depth fiber at \ensuremath{k}.
 it first sets \ensuremath{F} to \ensuremath{DepthFiber(m,\,k)}, and then
 chooses a permutation \ensuremath{P} of \ensuremath{F
 } such that \ensuremath{Delta\_\{P[r\.{+}1]\}} is contained in
 \ensuremath{Delta\_\{P[r]\}} for every
 successive pair.
 we do not impose any canonical ordering: \ensuremath{{\CHOOSE}} simply
 selects one
 admissible enumeration.
\end{cpar}%
\vshade{5.0}%
\begin{cpar}{0}{F}{F}{0}{0}{}%
\ensuremath{jIndex(M,\,k)
}%
\end{cpar}%
\begin{cpar}{0}{F}{F}{0}{0}{}%
 returns \ensuremath{j\_k}, \ensuremath{i.e}. the last index in the chosen
 admissible enumeration
 at depth \ensuremath{k}.
 if \ensuremath{k} is not a depth occurring in \ensuremath{M}, it returns 0.
 the last element is obtained using the general \ensuremath{Last} operator.
\end{cpar}%
\vshade{5.0}%
\begin{cpar}{0}{F}{F}{0}{0}{}%
\ensuremath{IVee(M)
}%
\end{cpar}%
\begin{cpar}{0}{F}{F}{0}{0}{}%
returns the permutation \ensuremath{i \.{\mapsto} i\.{\ct}vee}.
 for each depth, the chosen admissible enumeration is treated as a cycle.
 the set of these cycles is then passed to the general
 \ensuremath{PermutationFromCycles} operator.
\end{cpar}%
\vshade{5.0}%
\begin{cpar}{0}{F}{F}{0}{0}{}%
\ensuremath{SegmentTransform(M,\,i)
}%
\end{cpar}%
\begin{cpar}{0}{F}{F}{0}{0}{}%
returns \ensuremath{Delta}\mbox{'}\_i.
 its beginning is \ensuremath{b(Delta\_i)}, while its ending is
 \ensuremath{e(Delta\_\{i\.{\ct}vee\})}.
\end{cpar}%
\vshade{5.0}%
\begin{cpar}{0}{F}{F}{0}{0}{}%
\ensuremath{DistinguishedIndices(M)
}%
\end{cpar}%
\begin{cpar}{0}{F}{F}{0}{0}{}%
returns the set I\mbox{'} of distinguished indexes \ensuremath{j\_k},
 one for every depth occurring in \ensuremath{M}.
\end{cpar}%
\vshade{5.0}%
\begin{cpar}{0}{F}{F}{0}{0}{}%
\ensuremath{RemainingIndices(M)
}%
\end{cpar}%
\begin{cpar}{0}{F}{F}{0}{0}{}%
 returns I'' \ensuremath{\.{=}} I \ensuremath{\.{\,\backslash\,}} I\mbox{'},
 \ensuremath{i.e}. the global indexes that are not distinguished.
\end{cpar}%
\vshade{5.0}%
\begin{cpar}{0}{F}{F}{0}{0}{}%
\ensuremath{l(M)
}%
\end{cpar}%
\begin{cpar}{0}{F}{F}{0}{0}{}%
returns the ladder \ensuremath{l(m)}.
 for each depth \ensuremath{r\.{-}1} from 0 to \ensuremath{d(M)}, it takes
 the transformed segment
 Delta\mbox{'}\_\ensuremath{\{j\_\{r\.{-}1\}\}}.
\end{cpar}%
\vshade{5.0}%
\begin{cpar}{0}{F}{F}{0}{0}{}%
\ensuremath{DerivedMultisegment(M)
}%
\end{cpar}%
\begin{cpar}{0}{F}{F}{0}{0}{}%
 returns \ensuremath{m}\mbox{'}, the multisegment consisting of
 \ensuremath{Delta}\mbox{'}\_i for all \ensuremath{i} in I''.
 since a multisegment is represented by a sequence, we need some
 enumeration of these values.
 this is handled by the general \ensuremath{EnumerateValues} operator rather
 than being
 built directly into this paper-specific definition.
\end{cpar}%
\vshade{5.0}%
\begin{cpar}{0}{F}{F}{0}{0}{}%
\ensuremath{K(M)
}%
\end{cpar}%
\begin{cpar}{0}{F}{F}{0}{0}{}%
returns \ensuremath{{\langle}l(M),\, DerivedMultisegment(M){\rangle}},
 \ensuremath{i.e}. \ensuremath{K(m) \.{=}}
 (\ensuremath{l(m)},\ensuremath{m}\mbox{'}).
\end{cpar}%
\vshade{10.0}%
\begin{cpar}{0}{F}{F}{0}{0}{}%
\ensuremath{min(M)
}%
\end{cpar}%
\begin{cpar}{0}{F}{F}{0}{0}{}%
 returns \ensuremath{min} \ensuremath{m}, \ensuremath{i.e}. the smallest
 beginning among all segments of \ensuremath{M}.
 this is implemented by taking \ensuremath{Min} of the set of all
 \ensuremath{b(Delta\_i)}.
\end{cpar}%
\vshade{5.0}%
\begin{cpar}{0}{F}{F}{0}{0}{}%
\ensuremath{LeadingSequence(M)
}%
\end{cpar}%
\begin{cpar}{0}{F}{F}{0}{0}{}%
 chooses a sequence \ensuremath{{\langle}i\_1,\,\.{\dots},\,i\_k{\rangle}}
 satisfying \ensuremath{3.1a\.{-}3.1c}.
\end{cpar}%
\vshade{5.0}%
\begin{cpar}{0}{F}{F}{0}{0}{}%
 internally, \ensuremath{CandidateSequences} contains all possible nonempty
 sequences
 of indexes of length at most \ensuremath{Len(M)}.
\end{cpar}%
\vshade{5.0}%
\begin{cpar}{0}{F}{F}{0}{0}{}%
\ensuremath{FirstOK(i)} checks \ensuremath{3.1a}:
 \ensuremath{Delta\_i} begins at \ensuremath{min(M)}, and among all segments
 beginning there it has
 minimal ending.
\end{cpar}%
\vshade{5.0}%
\begin{cpar}{0}{F}{F}{0}{0}{}%
 \ensuremath{NextCandidate(i,\,q)} checks whether \ensuremath{q} can follow
 \ensuremath{i}:
 \ensuremath{Delta\_i} ≺ \ensuremath{Delta\_q} and
 \ensuremath{b(Delta\_q)\.{=}b(Delta\_i)\.{+}1}.
\end{cpar}%
\vshade{5.0}%
\begin{cpar}{0}{F}{F}{0}{0}{}%
\ensuremath{NextOK(i,\,next)} checks \ensuremath{3.1b}:
 next is a valid candidate, and among all valid candidates it has
 minimal ending.
\end{cpar}%
\vshade{5.0}%
\begin{cpar}{0}{F}{F}{0}{0}{}%
 \ensuremath{IsValid(S)} checks that the first element satisfies
 \ensuremath{FirstOK},
 every successive pair satisfies \ensuremath{NextOK}, and that there is no
 valid
 successor after the last element, which gives \ensuremath{3.1c}.
\end{cpar}%
\vshade{5.0}%
\begin{cpar}{0}{F}{F}{0}{0}{}%
\ensuremath{{\CHOOSE}} then selects one complete leading sequence.
 these auxiliary operators are local because they are only machinery
 for implementing the recursive construction described in the paper.
\end{cpar}%
\vshade{5.0}%
\begin{cpar}{0}{F}{F}{0}{0}{}%
\ensuremath{IStar(M)
}%
\end{cpar}%
\begin{cpar}{0}{F}{F}{0}{0}{}%
 returns I* \ensuremath{\.{=} \{i\_1,\,\.{\dots},\,i\_k\}}, the set of indexes
 occurring in the chosen
 leading sequence.
 this is simply \ensuremath{Range(LeadingSequence(M))}.
\end{cpar}%
\vshade{5.0}%
\begin{cpar}{0}{F}{F}{0}{0}{}%
\ensuremath{DeltaStar(M,\,i)
}%
\end{cpar}%
\begin{cpar}{0}{F}{F}{0}{0}{}%
returns \ensuremath{Delta\_i}*.
 if \ensuremath{i} is in I*, it applies \ensuremath{SegTruncateLeft} to
 \ensuremath{Delta\_i}.
 otherwise, it leaves \ensuremath{Delta\_i} unchanged.
\end{cpar}%
\vshade{5.0}%
\begin{cpar}{0}{F}{F}{0}{0}{}%
\ensuremath{MCross(M)
}%
\end{cpar}%
\begin{cpar}{0}{F}{F}{0}{0}{}%
returns \ensuremath{m\.{\ct}}cross.
 first it keeps exactly those indexes for which \ensuremath{Delta\_i}* is not
 \ensuremath{EmptySeg}.
 it then uses \ensuremath{EnumerateValues} to turn the surviving transformed
 segments
 into a sequence.
 therefore empty transformed segments are discarded, while repetitions
 of nonempty segments are preserved.
\end{cpar}%
\vshade{5.0}%
\begin{cpar}{0}{F}{F}{0}{0}{}%
\ensuremath{DeltaCircle(M)
}%
\end{cpar}%
\begin{cpar}{0}{F}{F}{0}{0}{}%
returns \ensuremath{Delta\.{\ct}}circle.
 if the chosen leading sequence has length \ensuremath{k}, then
 \ensuremath{Delta\.{\ct}}circle \ensuremath{\.{=}} [\ensuremath{min(M)},
 \ensuremath{min(M)\.{+}k\.{-}1}].
\end{cpar}%
\vshade{5.0}%
\begin{cpar}{0}{F}{F}{0}{0}{}%
\ensuremath{MW(M)
}%
\end{cpar}%
\begin{cpar}{0}{F}{F}{0}{0}{}%
returns \ensuremath{{\langle}MCross(M),\, DeltaCircle(M){\rangle}},
 \ensuremath{i.e}. \ensuremath{MW(m) \.{=}}
 (\ensuremath{m\.{\ct}}cross,\ensuremath{Delta\.{\ct}}circle).
\end{cpar}%
\vshade{10.0}%
\begin{cpar}{0}{F}{F}{0}{0}{}%
\ceqdash{76}
\end{cpar}%
\begin{cpar}{0}{F}{F}{0}{0}{}%
\ensuremath{Utils
}%
\end{cpar}%
\begin{cpar}{0}{F}{F}{0}{0}{}%
\ceqdash{76}
\end{cpar}%
\vshade{5.0}%
\begin{cpar}{0}{F}{F}{0}{0}{}%
\ensuremath{Range(f)
}%
\end{cpar}%
\begin{cpar}{0}{F}{F}{0}{0}{}%
returns the image/range of a function \ensuremath{f},
 \ensuremath{i.e}. \ensuremath{\{f[x] \.{:} x \.{\in} {\DOMAIN} f\}}.
\end{cpar}%
\vshade{5.0}%
\begin{cpar}{0}{F}{F}{0}{0}{}%
\ensuremath{Max(S)
}%
\end{cpar}%
\begin{cpar}{0}{F}{F}{0}{0}{}%
returns the maximum element of a nonempty finite set \ensuremath{S}.
\end{cpar}%
\vshade{5.0}%
\begin{cpar}{0}{F}{F}{0}{0}{}%
\ensuremath{Min(S)
}%
\end{cpar}%
\begin{cpar}{0}{F}{F}{0}{0}{}%
returns the minimum element of a nonempty finite set \ensuremath{S}.
\end{cpar}%
\vshade{5.0}%
\begin{cpar}{0}{F}{F}{0}{0}{}%
\ensuremath{IsInjective(f)
}%
\end{cpar}%
\begin{cpar}{0}{F}{F}{0}{0}{}%
checks whether \ensuremath{f} is injective:
 distinct elements of \ensuremath{{\DOMAIN} f} must have distinct images.
\end{cpar}%
\vshade{5.0}%
\begin{cpar}{0}{F}{F}{0}{0}{}%
\ensuremath{Permutations(S)
}%
\end{cpar}%
\begin{cpar}{0}{F}{F}{0}{0}{}%
returns the set of all permutations of a finite set \ensuremath{S}.
 a permutation is represented as a function from
 \ensuremath{1\.{\dotdot}Cardinality(S)} to \ensuremath{S
} which is injective.
 because the domain and \ensuremath{S} have the same finite cardinality,
 injectivity also guarantees that every element of \ensuremath{S} occurs
 exactly once.
\end{cpar}%
\vshade{5.0}%
\begin{cpar}{0}{F}{F}{0}{0}{}%
\ensuremath{ChoosePermutation(S)
}%
\end{cpar}%
\begin{cpar}{0}{F}{F}{0}{0}{}%
 chooses one arbitrary permutation of \ensuremath{S} using
 \ensuremath{{\CHOOSE}}.
 this is useful whenever we need to turn a set of indexes into a sequence
 but do not care about the particular order.
\end{cpar}%
\vshade{5.0}%
\begin{cpar}{0}{F}{F}{0}{0}{}%
\ensuremath{Fiber(f,\,y)
}%
\end{cpar}%
\begin{cpar}{0}{F}{F}{0}{0}{}%
returns the fiber/preimage of \ensuremath{y} under \ensuremath{f},
 \ensuremath{i.e}. all \ensuremath{x} in \ensuremath{{\DOMAIN} f} such that
 \ensuremath{f[x]\.{=}y}.
\end{cpar}%
\vshade{5.0}%
\begin{cpar}{0}{F}{F}{0}{0}{}%
\ensuremath{Last(S)
}%
\end{cpar}%
\begin{cpar}{0}{F}{F}{0}{0}{}%
returns the last element of a nonempty sequence \ensuremath{S},
 \ensuremath{i.e}. \ensuremath{S[Len(S)]}.
\end{cpar}%
\vshade{5.0}%
\begin{cpar}{0}{F}{F}{0}{0}{}%
\ensuremath{Multiplicity(S,\,x)
}%
\end{cpar}%
\begin{cpar}{0}{F}{F}{0}{0}{}%
 returns the number of indexes \ensuremath{i} for which
 \ensuremath{S[i]\.{=}x}.
 when \ensuremath{S} represents a bag/multiset, this is the multiplicity of
 \ensuremath{x
} in that bag.
\end{cpar}%
\vshade{5.0}%
\begin{cpar}{0}{F}{F}{0}{0}{}%
\ensuremath{SubBag(A,\,B)
}%
\end{cpar}%
\begin{cpar}{0}{F}{F}{0}{0}{}%
checks whether A is a subbag of \ensuremath{B}.
 for every value occurring in A, its multiplicity in A must be no greater
 than its multiplicity in \ensuremath{B}.
\end{cpar}%
\vshade{5.0}%
\begin{cpar}{0}{F}{F}{0}{0}{}%
\ensuremath{SameBag(A,\,B)
}%
\end{cpar}%
\begin{cpar}{0}{F}{F}{0}{0}{}%
checks whether A and \ensuremath{B} represent the same bag.
 this means \ensuremath{SubBag(A,\,B)} and \ensuremath{SubBag(B,\,A)}, so
 order is ignored but
 repetitions/multiplicities are preserved.
\end{cpar}%
\vshade{5.0}%
\begin{cpar}{0}{F}{F}{0}{0}{}%
\ensuremath{PermutationFromCycles(I,\,Cycles)
}%
\end{cpar}%
\begin{cpar}{0}{F}{F}{0}{0}{}%
constructs a permutation of I from a set of cycles.
 for each \ensuremath{i}, it finds the cycle \ensuremath{E} containing
 \ensuremath{i} and the position \ensuremath{r} of \ensuremath{i} in
 \ensuremath{E}.
 if \ensuremath{i} is the last element of \ensuremath{E}, it maps
 \ensuremath{i} to \ensuremath{E[1]}.
 otherwise, it maps \ensuremath{i} to \ensuremath{E[r\.{+}1]}.
\end{cpar}%
\vshade{5.0}%
\begin{cpar}{0}{F}{F}{0}{0}{}%
\ensuremath{EnumerateValues(f)
}%
\end{cpar}%
\begin{cpar}{0}{F}{F}{0}{0}{}%
turns the values of a finite function \ensuremath{f} into a sequence.
 it chooses a permutation of \ensuremath{{\DOMAIN} f} and returns the values
 of \ensuremath{f} in
 that order.
 unlike \ensuremath{Range(f)}, this preserves repetitions:
 if two different indexes map to the same value, that value appears twice
 in the resulting sequence.
\end{cpar}%
\vshade{10.0}%
\begin{cpar}{0}{F}{F}{0}{0}{}%
\ceqdash{76}
\end{cpar}%
\begin{cpar}{0}{F}{F}{0}{0}{}%
Notes on notation and organization
\end{cpar}%
\begin{cpar}{0}{F}{F}{0}{0}{}%
\ceqdash{76}
\end{cpar}%
\vshade{5.0}%
\begin{cpar}{0}{F}{F}{0}{0}{}%
why choose operators and not function notations\.{?}
 for instance, why define Multiplicity as \ensuremath{Multiplicity(S,\,x)}
 without
 specifying the domain, meaning something like
 \ensuremath{Multiplicity[S \.{\in} Seq(\.{\dots}),\, x \.{\in}
 \.{\dots}]}\.{?}
\end{cpar}%
\vshade{5.0}%
\begin{cpar}{0}{F}{F}{0}{0}{}%
TLA+ is untyped, so if we do need to specify the domain, we can do so in
 the definitions that use these objects.
 second, the bracket notation defines a function, while the operator
 notation defines an operator.
 from a more practical standpoint, unfolding the operator notation in a
 proof requires less machinery and less proof obligations.
\end{cpar}%
\vshade{15.0}%
\begin{cpar}{0}{F}{F}{0}{0}{}%
 the purpose of \ensuremath{Utils} is to separate operations that are not
 specific to
 this paper, such as \ensuremath{Range}, \ensuremath{Fiber}, permutations,
 bag equality, or
 constructing a permutation from cycles.
\end{cpar}%
\end{lcom}%
\end{document}
